\documentclass[11pt]{article}
\usepackage[margin=1in]{geometry}
\usepackage{booktabs}
\usepackage{longtable}
\usepackage{array}
\usepackage{amsmath}
\usepackage[section]{placeins}
\usepackage{hyperref}
\hypersetup{colorlinks=true, linkcolor=blue, urlcolor=blue}
\sloppy
\newcolumntype{L}[1]{>{\raggedright\arraybackslash}p{#1}}

% Every number in this document is read from the memo or artefact that owns it: the memos'
% own generated macros files are input below, and build_chronicle.py writes the rest
% (tables/chronicle_macros.tex) from the older memos' JSON artefacts, git history and the data.
% GENERATED BY build_chronicle.py --- do not edit by hand.
\newcommand{\chClockFirst}{1985}
\newcommand{\chClockLast}{2026}
\newcommand{\chCondFirst}{1984}
\newcommand{\chCondLast}{2026}
\newcommand{\chDateClocks}{2026-09-11}
\newcommand{\chDateConditionsContent}{2026-09-11}
\newcommand{\chDateConditionsOfApproval}{2026-09-08}
\newcommand{\chDateDataAcquisition}{2026-09-08}
\newcommand{\chDateDatasfPlanningCatalogue}{2026-09-11}
\newcommand{\chDateDiscretionaryReviewPatterns}{2026-09-07}
\newcommand{\chDateExactions}{2026-09-11}
\newcommand{\chDateExtractionMethodComparison}{2026-09-07}
\newcommand{\chDateMeetingLevelInfo}{2026-09-04}
\newcommand{\chDatePermitLinkage}{2026-09-08}
\newcommand{\chDatePermitsContent}{2026-09-11}
\newcommand{\chDatePredictingDelay}{2026-09-08}
\newcommand{\chDateRegulatoryTimeline}{2026-09-11}
\newcommand{\chDbiCondTalk}{87}
\newcommand{\chExFirst}{2001}
\newcommand{\chExLast}{2026}
\newcommand{\chGoldItems}{294}
\newcommand{\chHaikuTest}{96.6}
\newcommand{\chItemFirst}{1998}
\newcommand{\chItemLast}{2026}
\newcommand{\chItems}{16{,}199}
\newcommand{\chLawFirst}{2016}
\newcommand{\chLawLast}{2025}
\newcommand{\chLaws}{45}
\newcommand{\chMeetAcc}{93.9}
\newcommand{\chMeetEarly}{91.0}
\newcommand{\chMeetFirst}{1998}
\newcommand{\chMeetLast}{2026}
\newcommand{\chMeetLate}{98.3}
\newcommand{\chMeetScored}{134}
\newcommand{\chMeetValues}{938}
\newcommand{\chMeetings}{1{,}219}
\newcommand{\chMemos}{13}
\newcommand{\chPagesClocks}{17}
\newcommand{\chPagesConditionsContent}{66}
\newcommand{\chPagesConditionsOfApproval}{13}
\newcommand{\chPagesDataAcquisition}{36}
\newcommand{\chPagesDatasfPlanningCatalogue}{43}
\newcommand{\chPagesDiscretionaryReviewPatterns}{17}
\newcommand{\chPagesExactions}{19}
\newcommand{\chPagesExtractionMethodComparison}{16}
\newcommand{\chPagesMeetingLevelInfo}{17}
\newcommand{\chPagesPermitLinkage}{9}
\newcommand{\chPagesPermitsContent}{21}
\newcommand{\chPagesPredictingDelay}{10}
\newcommand{\chPagesRegulatoryTimeline}{20}
\newcommand{\chPagesTotal}{304}
\newcommand{\chParcelPrecision}{0.72}
\newcommand{\chParcelRecall}{0.92}
\newcommand{\chPermitMatch}{97.6}
\newcommand{\chPermitsDistinct}{2{,}295}
\newcommand{\chPriceFirst}{1998}
\newcommand{\chPriceLast}{2026}
\newcommand{\chRecFirst}{1920}
\newcommand{\chRecLast}{2026}
\newcommand{\chRegexTest}{75.2}
\newcommand{\chRowsClocks}{34{,}320}
\newcommand{\chRowsConditions}{94{,}169}
\newcommand{\chRowsEnvelope}{5{,}824{,}609}
\newcommand{\chRowsExactions}{3{,}574}
\newcommand{\chRowsNumeric}{6{,}970}
\newcommand{\chRowsOutcomes}{10{,}253}
\newcommand{\chRrFirst}{1998}
\newcommand{\chRrLast}{2009}
\newcommand{\chRrWith}{1{,}032}
\newcommand{\chRrWithout}{554}

% GENERATED BY analyze_corpus.py universe --- do not edit by hand.
\newcommand{\corpusCaseMerges}{29}
\newcommand{\corpusCasesPrinted}{8{,}987}
\newcommand{\corpusMaxDaysPrinted}{3{,}423}
\newcommand{\corpusMaxDaysUniverse}{3{,}745}
\newcommand{\corpusRepeatPrinted}{37.9}
\newcommand{\corpusRepeatUniverse}{38.1}
\newcommand{\corpusUniverse}{8{,}958}

% GENERATED BY analyze_conditions.py report --- do not edit by hand.
\newcommand{\condBridgeCUItems}{2{,}348}
\newcommand{\condBridgeCUPct}{46.3}
\newcommand{\condCUItems}{5{,}074}
\newcommand{\condCaseMerges}{29}
\newcommand{\condCasesPrinted}{8{,}987}
\newcommand{\condCondItems}{4{,}048}
\newcommand{\condCondTextCU}{2{,}522}
\newcommand{\condCondWithText}{1{,}952}
\newcommand{\condCondWithTextPct}{48.2}
\newcommand{\condEarliestAny}{1984}
\newcommand{\condEarliestFtp}{2009}
\newcommand{\condEarliestInUniverse}{1999}
\newcommand{\condEarliestItemYear}{1999}
\newcommand{\condEarliestModules}{2007}
\newcommand{\condEarliestWayback}{2001}
\newcommand{\condOverlapCU}{1{,}424}
\newcommand{\condProbeHits}{10{,}165}
\newcommand{\condProbed}{127{,}920}
\newcommand{\condReachEarlyHtml}{2.4}
\newcommand{\condReachEarlyPdf}{78.9}
\newcommand{\condReachLateHtml}{7.6}
\newcommand{\condResPreTen}{1{,}600}
\newcommand{\condResPreTenNone}{1{,}586}
\newcommand{\condRowsAll}{94{,}169}
\newcommand{\condSectionsWithText}{4{,}180}
\newcommand{\condUniverse}{8{,}958}

% GENERATED BY analyze_delay.py universe --- do not edit by hand.
\newcommand{\delayCasesPrinted}{8{,}987}
\newcommand{\delayCasesUniverse}{8{,}958}
\newcommand{\delayMultiPrinted}{39.1}
\newcommand{\delayMultiUniverse}{39.2}
\newcommand{\delayPositivePrinted}{34.1}
\newcommand{\delayPositiveUniverse}{34.3}
\newcommand{\delaySamplePrinted}{8{,}382}
\newcommand{\delaySampleUniverse}{8{,}353}

% GENERATED BY acquire_external_data.py — do not edit by hand.
% Every number that appears in the memo's prose is defined here.
\newcommand{\acqAsrBalanced}{4{,}372}
\newcommand{\acqAsrFirst}{2007}
\newcommand{\acqAsrFound}{5{,}565}
\newcommand{\acqAsrGapYears}{9}
\newcommand{\acqAsrJoin}{81.6}
\newcommand{\acqAsrLast}{2025}
\newcommand{\acqAsrNYears}{19}
\newcommand{\acqAsrParcels}{222{,}251}
\newcommand{\acqAsrParcelsPerYear}{208{,}927}
\newcommand{\acqAsrRows}{3{,}934{,}467}
\newcommand{\acqAsrSought}{8{,}294}
\newcommand{\acqBridgeCU}{46.0}
\newcommand{\acqBridgeCUN}{2{,}336}
\newcommand{\acqBridgeCUTotal}{5{,}074}
\newcommand{\acqBridgeDR}{71.8}
\newcommand{\acqBridgeDbi}{7{,}833}
\newcommand{\acqBridgeDbiShare}{48.4}
\newcommand{\acqBridgeDnt}{57.3}
\newcommand{\acqBridgeItems}{7{,}875}
\newcommand{\acqBridgeLpa}{70.5}
\newcommand{\acqBridgePca}{7.6}
\newcommand{\acqBridgeShare}{48.6}
\newcommand{\acqBridgeVar}{63.3}
\newcommand{\acqCaseAny}{92.7}
\newcommand{\acqCaseAnyN}{8{,}334}
\newcommand{\acqCaseJoinBreakYear}{2002}
\newcommand{\acqCaseJoinLateMin}{96}
\newcommand{\acqCaseNorm}{75.0}
\newcommand{\acqCaseRaw}{74.7}
\newcommand{\acqCaseStem}{92.7}
\newcommand{\acqCases}{8{,}987}
\newcommand{\acqClDenseFirst}{1968}
\newcommand{\acqClFirst}{1900}
\newcommand{\acqClItemKeys}{8{,}294}
\newcommand{\acqClItemKeysTraded}{3{,}882}
\newcommand{\acqClItemKeysTradedPct}{46.8}
\newcommand{\acqClItemsEver}{7{,}727}
\newcommand{\acqClItemsEverPct}{56.7}
\newcommand{\acqClItemsNear}{3{,}125}
\newcommand{\acqClItemsNearPct}{22.9}
\newcommand{\acqClKeyed}{100.0}
\newcommand{\acqClLast}{2024}
\newcommand{\acqClMaxDate}{2024-05-16}
\newcommand{\acqClMedianRecent}{1.42}
\newcommand{\acqClParcels}{174{,}102}
\newcommand{\acqClPriced}{415{,}854}
\newcommand{\acqClResidential}{320{,}323}
\newcommand{\acqClRows}{447{,}883}
\newcommand{\acqClWindowYears}{3}
\newcommand{\acqCoreLogicSales}{415{,}854}
\newcommand{\acqCorpusSpan}{1998--2026}
\newcommand{\acqCuHousingProd}{17.6}
\newcommand{\acqCuMohcd}{1.5}
\newcommand{\acqCuPipeline}{4.5}
\newcommand{\acqCuUnits}{91.9}
\newcommand{\acqCuUnitsNZ}{27.1}
\newcommand{\acqCuUnitsNZearly}{18.3}
\newcommand{\acqCuUnitsNZlate}{30.9}
\newcommand{\acqDelayIssueMedian}{245}
\newcommand{\acqDelayIssuePninety}{768}
\newcommand{\acqDemandFiles}{76}
\newcommand{\acqDemandGB}{4.92}
\newcommand{\acqDistinctParcelKeys}{8{,}294}
\newcommand{\acqEitherCU}{46.1}
\newcommand{\acqEitherDR}{93.2}
\newcommand{\acqEitherRoute}{54.0}
\newcommand{\acqFeeAreaRows}{33}
\newcommand{\acqFeeEarliestSnapshot}{2019-09-15}
\newcommand{\acqFeeFirst}{2019}
\newcommand{\acqFeeLast}{2026}
\newcommand{\acqFeeRates}{1{,}341}
\newcommand{\acqFeeRegisters}{8}
\newcommand{\acqInclusionaryFirst}{199.50}
\newcommand{\acqInclusionaryLast}{249.66}
\newcommand{\acqItems}{16{,}199}
\newcommand{\acqLagCU}{175}
\newcommand{\acqLagCoverage}{92.7}
\newcommand{\acqLagDR}{105}
\newcommand{\acqLagMedian}{142}
\newcommand{\acqLagN}{8{,}334}
\newcommand{\acqLagNegative}{0.3}
\newcommand{\acqLagPninety}{758}
\newcommand{\acqLagPninetynine}{2346}
\newcommand{\acqLagTrunc}{1{,}500}
\newcommand{\acqLagVar}{280}
\newcommand{\acqLandUseMulti}{77}
\newcommand{\acqMinCasesFig}{10}
\newcommand{\acqMinItemsBridge}{100}
\newcommand{\acqModernUnmatched}{92}
\newcommand{\acqMohcdMaxYear}{2030}
\newcommand{\acqMohcdRows}{190}
\newcommand{\acqNonProjRows}{233{,}394}
\newcommand{\acqNonUniqueSources}{14}
\newcommand{\acqNpApplicant}{66.9}
\newcommand{\acqNpApplicantOrg}{34.3}
\newcommand{\acqNpBuildingPermits}{22.9}
\newcommand{\acqNpCloseDate}{94.9}
\newcommand{\acqNpMinYear}{1920}
\newcommand{\acqNpOpenDate}{100.0}
\newcommand{\acqNpParentId}{26.7}
\newcommand{\acqNpPlanner}{77.7}
\newcommand{\acqNpRecordStatus}{100.0}
\newcommand{\acqOldFormatCases}{536}
\newcommand{\acqOldFormatMatched}{0}
\newcommand{\acqParcelActive}{66.9}
\newcommand{\acqParcelJoin}{89.5}
\newcommand{\acqParcelJoinN}{12{,}195}
\newcommand{\acqParcelNoBlock}{193}
\newcommand{\acqParcelNoLot}{1{,}232}
\newcommand{\acqParcelRetired}{3{,}080}
\newcommand{\acqParcelRows}{236{,}556}
\newcommand{\acqParcelShare}{84.1}
\newcommand{\acqPipelineRows}{1{,}917}
\newcommand{\acqPrBuildingPermits}{45.2}
\newcommand{\acqPrDecision}{18.3}
\newcommand{\acqPrDecisionDate}{17.3}
\newcommand{\acqPrEnvDoc}{25.1}
\newcommand{\acqPrInclFlag}{2.9}
\newcommand{\acqPrUnitsAff}{100.0}
\newcommand{\acqPrUnitsNet}{94.0}
\newcommand{\acqPrUnitsProp}{100.0}
\newcommand{\acqPrintedAll}{22.1}
\newcommand{\acqPrintedCU}{0.3}
\newcommand{\acqPrintedDR}{85.5}
\newcommand{\acqProjRows}{54{,}515}
\newcommand{\acqRetrieved}{2026-09-08}
\newcommand{\acqScaleEarliest}{2005}
\newcommand{\acqSourcesCached}{19}
\newcommand{\acqTsfNinetyNineFirst}{10.29}
\newcommand{\acqTsfNinetyNineLast}{13.39}
\newcommand{\acqUcMaxYear}{2205}
\newcommand{\acqWithParcel}{13{,}620}
\newcommand{\acqZhviFirst}{2000-01}
\newcommand{\acqZhviLast}{2026-07}
\newcommand{\acqZhviMonths}{319}
\newcommand{\acqZhviZips}{26}
\newcommand{\acqZoningParcelSpan}{1998--2008}
\newcommand{\acqZoningParcelYears}{9}
\newcommand{\acqZoningPolygonYears}{11}
\newcommand{\acqZoriFirst}{2015-01}
\newcommand{\acqZoriLast}{2026-07}
\newcommand{\acqZoriMonths}{139}
\newcommand{\acqZoriZips}{24}

% GENERATED BY acquire_external_data.py envelope --- do not edit by hand.
\newcommand{\envClaims}{25}
\newcommand{\envClaimsOk}{25}
\newcommand{\envDblMin}{5}
\newcommand{\envDblPct}{50}
\newcommand{\envDblTotal}{299{,}253}
\newcommand{\envDensityBinds}{90}
\newcommand{\envFloor}{10}
\newcommand{\envGpu}{841}
\newcommand{\envHeight}{99}
\newcommand{\envKindNone}{4}
\newcommand{\envKindNp}{1}
\newcommand{\envKindPerArea}{2}
\newcommand{\envKindPerAreaThree}{8}
\newcommand{\envKindPerLot}{78}
\newcommand{\envKindUnread}{7}
\newcommand{\envLot}{94}
\newcommand{\envParcels}{155{,}367}
\newcommand{\envPriced}{91}
\newcommand{\envRead}{93}
\newcommand{\envReadFrom}{2020}
\newcommand{\envReadTo}{2021}
\newcommand{\envSensHi}{937{,}959}
\newcommand{\envSensLo}{687{,}630}
\newcommand{\envUnitsEight}{694{,}338}
\newcommand{\envUnitsSixteen}{781{,}591}
\newcommand{\envUnitsTotal}{795{,}840}
\newcommand{\envYearA}{2008}
\newcommand{\envYearB}{2016}
\newcommand{\envYearLast}{2026}

% GENERATED BY analyze_datasf_catalogue.py --- do not edit by hand.
\newcommand{\catApiTotal}{273}
\newcommand{\catCached}{36}
\newcommand{\catCharts}{1}
\newcommand{\catCore}{35}
\newcommand{\catDatasets}{155}
\newcommand{\catDepartments}{23}
\newcommand{\catDomainsOnly}{2}
\newcommand{\catDomainsOnlyOld}{2}
\newcommand{\catFiles}{2}
\newcommand{\catFilters}{22}
\newcommand{\catGeom}{130}
\newcommand{\catKeyAny}{61}
\newcommand{\catKeyBlklot}{34}
\newcommand{\catKeyBlockLot}{30}
\newcommand{\catKeyCase}{18}
\newcommand{\catKeyPermit}{20}
\newcommand{\catLinks}{13}
\newcommand{\catMaps}{78}
\newcommand{\catNonDatasets}{118}
\newcommand{\catNotRelevant}{155}
\newcommand{\catOldHostCount}{273}
\newcommand{\catOldHostSameIds}{the same}
\newcommand{\catPlanningDept}{158}
\newcommand{\catPossible}{83}
\newcommand{\catPossibleNotCached}{82}
\newcommand{\catRetrieved}{2026-09-11}
\newcommand{\catRowsTotal}{28{,}161{,}918}
\newcommand{\catStale}{174}
\newcommand{\catStaleCutoff}{2024-09-11}
\newcommand{\catStaleYears}{2}
\newcommand{\catStories}{2}
\newcommand{\catTotal}{273}
\newcommand{\catUnassessed}{0}
\newcommand{\catWithRows}{255}
\newcommand{\catInvCommission}{1}
\newcommand{\catInvPipelineInternal}{4}
\newcommand{\catInvPlanning}{109}
\newcommand{\catInvPublished}{89}
\newcommand{\catInvUnpublished}{20}
\newcommand{\catRelPagingDistinct}{261}
\newcommand{\catRelPagingRows}{273}

% GENERATED BY collect_conditions.py report --- do not edit by hand.
\newcommand{\ccAdoptedWithText}{1{,}794}
\newcommand{\ccAnatAdopted}{2013-12-19}
\newcommand{\ccAnatCase}{2013.1613C}
\newcommand{\ccAnatConditions}{12}
\newcommand{\ccAnatMotion}{19048}
\newcommand{\ccAnatPages}{13}
\newcommand{\ccAnatUrl}{https://commissions.sfplanning.org/cpcmotions/2013/19048.pdf}
\newcommand{\ccCandidates}{127{,}920}
\newcommand{\ccCaseMerges}{29}
\newcommand{\ccCasesPrinted}{8{,}987}
\newcommand{\ccCdxFirstMotion}{2002}
\newcommand{\ccCdxFirstPacket}{2011}
\newcommand{\ccCdxFirstYear}{1999}
\newcommand{\ccCdxMotionsHtml}{20}
\newcommand{\ccCdxMotionsPdf}{136}
\newcommand{\ccCdxOk}{92{,}055}
\newcommand{\ccCdxPackets}{2{,}963}
\newcommand{\ccCdxPdf}{6{,}088}
\newcommand{\ccCdxPrefixes}{11}
\newcommand{\ccCdxUrls}{97{,}503}
\newcommand{\ccCheckN}{30}
\newcommand{\ccCondItems}{4{,}048}
\newcommand{\ccCondNoText}{2{,}098}
\newcommand{\ccCondNoTextNoNumber}{620}
\newcommand{\ccCondNoTextWithNumber}{1{,}478}
\newcommand{\ccCondWithText}{1{,}950}
\newcommand{\ccCondWithTextPct}{48.2}
\newcommand{\ccConditionRows}{84{,}684}
\newcommand{\ccCountVerdict}{is not rising with}
\newcommand{\ccDiffAdded}{1{,}148}
\newcommand{\ccDiffAnyChange}{467}
\newcommand{\ccDiffAnyChangePct}{36}
\newcommand{\ccDiffByText}{72}
\newcommand{\ccDiffCarriedPct}{94}
\newcommand{\ccDiffChanged}{394}
\newcommand{\ccDiffDraftConds}{19{,}749}
\newcommand{\ccDiffIdentical}{17{,}133}
\newcommand{\ccDiffMods}{224}
\newcommand{\ccDiffModsChange}{183}
\newcommand{\ccDiffModsNoChange}{41}
\newcommand{\ccDiffNear}{1{,}411}
\newcommand{\ccDiffNoItem}{45}
\newcommand{\ccDiffNoModsChange}{257}
\newcommand{\ccDiffNoModsNoChange}{755}
\newcommand{\ccDiffPairs}{1{,}281}
\newcommand{\ccDiffRemoved}{811}
\newcommand{\ccDiffSame}{90}
\newcommand{\ccDiffUnchanged}{18{,}544}
\newcommand{\ccDiscoveredCensus}{8{,}932}
\newcommand{\ccDiscoveredUrls}{11{,}527}
\newcommand{\ccDocsRead}{9{,}956}
\newcommand{\ccDraftWithText}{2{,}376}
\newcommand{\ccDriftN}{40}
\newcommand{\ccDriftRewriteMax}{0.60}
\newcommand{\ccDriftRewritten}{Additional Project Authorization; Eastern Neighborhoods Infrastructure Impact Fee; Lighting; Mitigation Measures; Noise; Extension; Hours of Operation}
\newcommand{\ccDriftRewrittenN}{7}
\newcommand{\ccDriftSame}{27}
\newcommand{\ccDriftSameMin}{0.99}
\newcommand{\ccEarliestAdopted}{1984}
\newcommand{\ccEarliestAny}{1984}
\newcommand{\ccEarliestDraft}{2007}
\newcommand{\ccEarliestEmbedded}{1984}
\newcommand{\ccEarliestFtp}{2009}
\newcommand{\ccEarliestInUniverse}{1994}
\newcommand{\ccEarliestItemYear}{1999}
\newcommand{\ccEarliestModules}{2007}
\newcommand{\ccEarliestPackets}{2007}
\newcommand{\ccEarliestVault}{2017}
\newcommand{\ccEarliestWayback}{2001}
\newcommand{\ccExtensionN}{3{,}366}
\newcommand{\ccExtensionWithPeriod}{30}
\newcommand{\ccFirstCitypln}{2008}
\newcommand{\ccFirstCpcdraFtp}{2003}
\newcommand{\ccFirstCpcdraMirror}{2011}
\newcommand{\ccFirstCpcmotionsFtp}{2010}
\newcommand{\ccFirstCpcmotionsYear}{2011}
\newcommand{\ccFirstCpcmotionsYearMirror}{2011}
\newcommand{\ccFirstCpcpackets}{2001}
\newcommand{\ccFirstCpcpacketsFtp}{2009}
\newcommand{\ccFirstCpcpacketsMirror}{2003}
\newcommand{\ccFirstMinutesVault}{2015}
\newcommand{\ccFirstOldsiteModules}{2008}
\newcommand{\ccGoldN}{20}
\newcommand{\ccHeadedRows}{57{,}324}
\newcommand{\ccHeadingsDistinct}{1{,}212}
\newcommand{\ccHeadingsOnce}{516}
\newcommand{\ccHeadingsTopShare}{76.9}
\newcommand{\ccHearingDocs}{4{,}007}
\newcommand{\ccHearingPages}{315}
\newcommand{\ccHits}{10{,}165}
\newcommand{\ccHtmlLinksFrom}{2000}
\newcommand{\ccHtmlLinksTo}{2017}
\newcommand{\ccKindBespoke}{Signage; Parking Maximum; Additional Project Authorization; Mitigation Measures; Screening - WTS}
\newcommand{\ccKindBespokeMax}{50}
\newcommand{\ccKindBespokeN}{5}
\newcommand{\ccKindBoiler}{Enforcement; Revocation due to Violation of Conditions; Conformity with Current Law; Validity; Expiration and Renewal; Diligent Pursuit; Sidewalk Maintenance; Final Materials; Rooftop Mechanical Equipment; Managing Traffic During Construction; Lighting; Garbage, Recycling, and Composting Receptacles; First Source Hiring; Transportation Sustainability Fee; Residential Child Care Impact Fee; Streetscape Plan; Transportation Demand Management (TDM) Program; Noise Control; Anti-Discriminatory Housing; Overhead Wiring; Compatibility with City Emergency Services -- WTS; Installation - WTS; Notification prior to Project Implementation Report - WTS; Implementation and Monitoring - WTS; Noise and Heat -- WTS}
\newcommand{\ccKindBoilerMin}{80}
\newcommand{\ccKindBoilerN}{25}
\newcommand{\ccKindNumber}{Bicycle Parking; Validity and Expiration; Hours of Operation}
\newcommand{\ccKindNumberMin}{75}
\newcommand{\ccKindNumberN}{3}
\newcommand{\ccListingsIndex}{0}
\newcommand{\ccListingsTried}{10}
\newcommand{\ccMinMotionsYear}{10}
\newcommand{\ccMinutesLinkDocs}{650}
\newcommand{\ccMinutesLinks}{12{,}917}
\newcommand{\ccOcrPages}{4{,}851}
\newcommand{\ccOldFound}{225}
\newcommand{\ccOldPerYear}{20}
\newcommand{\ccOldProbed}{578}
\newcommand{\ccOldsiteLinks}{2{,}127}
\newcommand{\ccOldsitePages}{247}
\newcommand{\ccOptionRows}{12{,}817}
\newcommand{\ccPerMotionFirstMed}{11}
\newcommand{\ccPerMotionFirstYear}{2001}
\newcommand{\ccPerMotionLastMed}{10}
\newcommand{\ccPerMotionLastYear}{2026}
\newcommand{\ccPerMotionMean}{15.5}
\newcommand{\ccPerMotionMedian}{14}
\newcommand{\ccPerMotionRho}{-0.26}
\newcommand{\ccPerMotionYears}{18}
\newcommand{\ccPerUnitsLargeMed}{26}
\newcommand{\ccPerUnitsLargeN}{338}
\newcommand{\ccPerUnitsNoneMed}{12}
\newcommand{\ccPreambleBlocks}{5}
\newcommand{\ccProbeFinished}{finished}
\newcommand{\ccProbeHits}{10{,}165}
\newcommand{\ccProbeLogin}{275}
\newcommand{\ccProbeMissing}{117{,}472}
\newcommand{\ccProbeOther}{0}
\newcommand{\ccProbeSoft}{8}
\newcommand{\ccProbed}{127{,}920}
\newcommand{\ccQuoteWords}{80}
\newcommand{\ccRatePerHost}{2}
\newcommand{\ccReachEarlyHtml}{2.4}
\newcommand{\ccReachEarlyPdf}{78.9}
\newcommand{\ccReachLateHtml}{7.5}
\newcommand{\ccReachMidPdf}{63.8}
\newcommand{\ccReachRecent}{84.0}
\newcommand{\ccReachTwentyTen}{71.6}
\newcommand{\ccReachVault}{83.0}
\newcommand{\ccResFoundRead}{381}
\newcommand{\ccResFoundUnread}{15}
\newcommand{\ccResNoneNoNumber}{583}
\newcommand{\ccResNoneNumber}{1{,}119}
\newcommand{\ccResPreTen}{1{,}601}
\newcommand{\ccResPreTenNone}{1{,}587}
\newcommand{\ccRowsAll}{91{,}720}
\newcommand{\ccSecLocated}{4{,}173}
\newcommand{\ccSecLocatedNone}{3}
\newcommand{\ccSecMention}{4{,}572}
\newcommand{\ccSecMentionNotLocated}{326}
\newcommand{\ccSecRead}{5{,}826}
\newcommand{\ccSecRejected}{90}
\newcommand{\ccSectionsOutsideUniverse}{448}
\newcommand{\ccSectionsWithText}{4{,}170}
\newcommand{\ccStrictRows}{62{,}239}
\newcommand{\ccTailN}{20}
\newcommand{\ccTaxCategories}{13}
\newcommand{\ccTaxDistinct}{53.5}
\newcommand{\ccTaxHeadings}{1{,}212}
\newcommand{\ccTaxObligation}{65.5}
\newcommand{\ccTaxProject}{16.4}
\newcommand{\ccTaxUnheaded}{85.4}
\newcommand{\ccTaxWeighted}{87.8}
\newcommand{\ccTopHeadings}{40}
\newcommand{\ccUncoveredTop}{First Source Hiring; Streetscape Plan; Anti-Discriminatory Housing; Overhead Wiring; Eating and Drinking Uses; Transformer Vault}
\newcommand{\ccUnheadedRows}{4{,}915}
\newcommand{\ccUniverse}{8{,}958}
\newcommand{\ccValAdjudicated}{1}
\newcommand{\ccValFrozen}{2026-09-11 03:07}
\newcommand{\ccValGoldConds}{395}
\newcommand{\ccValGoldDocs}{20}
\newcommand{\ccValInHead}{100.0}
\newcommand{\ccValInHeadN}{239}
\newcommand{\ccValInPrec}{97.7}
\newcommand{\ccValInRecall}{98.0}
\newcommand{\ccValNoneDocs}{2}
\newcommand{\ccValNoneFalse}{0}
\newcommand{\ccValOutDocs}{30}
\newcommand{\ccValOutEnd}{96.7}
\newcommand{\ccValOutHead}{96.6}
\newcommand{\ccValOutHeadN}{29}
\newcommand{\ccValOutPrec}{96.7}
\newcommand{\ccValOutRows}{30}
\newcommand{\ccValOutStart}{96.7}
\newcommand{\ccValRoneFrozen}{2026-09-11 02:30}
\newcommand{\ccValRoneInPrec}{85.1}
\newcommand{\ccValRoneInRecall}{76.5}
\newcommand{\ccValidityModal}{36}
\newcommand{\ccValidityModalShare}{96}
\newcommand{\ccValidityN}{3{,}904}
\newcommand{\ccValidityStatePct}{98}
\newcommand{\ccValidityYears}{18}
\newcommand{\ccVaultLinksFrom}{2017}

% GENERATED BY analyze_permit_content.py --- do not edit by hand.
\newcommand{\pcCellDaysMedB}{303}
\newcommand{\pcCellMedGap}{197}
\newcommand{\pcCols}{57}
\newcommand{\pcCompleteNoDate}{0.00}
\newcommand{\pcCostNominal}{1}
\newcommand{\pcCostPlaceholder}{10.2}
\newcommand{\pcCostPlaceholderLinked}{1.8}
\newcommand{\pcCostPlaceholderUnlinked}{2.0}
\newcommand{\pcCostQ}{5}
\newcommand{\pcCostQAltHi}{10}
\newcommand{\pcCostQAltLo}{2}
\newcommand{\pcDataCols}{53}
\newcommand{\pcDrBothN}{2{,}269}
\newcommand{\pcDrCases}{2{,}524}
\newcommand{\pcDrDbiMed}{294}
\newcommand{\pcDrDbiN}{2{,}356}
\newcommand{\pcDrDbiNeg}{3.0}
\newcommand{\pcDrDbiPninety}{791}
\newcommand{\pcDrDiffMed}{150}
\newcommand{\pcDrFiled}{2006--2020}
\newcommand{\pcDrNamedInDbi}{487}
\newcommand{\pcDrNegN}{201}
\newcommand{\pcDrNegOffParcel}{18}
\newcommand{\pcDrNegOwn}{110}
\newcommand{\pcDrNegPct}{8.9}
\newcommand{\pcDrNegRecord}{76}
\newcommand{\pcDrNegStem}{15}
\newcommand{\pcDrNegStemOff}{2}
\newcommand{\pcDrRecMed}{101}
\newcommand{\pcDrRecN}{2{,}418}
\newcommand{\pcDrRecPninety}{315}
\newcommand{\pcDrpN}{693}
\newcommand{\pcDrpNamed}{646}
\newcommand{\pcDrpYears}{2015--2020}
\newcommand{\pcDupIds}{0}
\newcommand{\pcFcdFilled}{1.2}
\newcommand{\pcFieldsDiffer}{0.01}
\newcommand{\pcFiledFirst}{1901}
\newcommand{\pcFiledLast}{2026}
\newcommand{\pcInclusionary}{10}
\newcommand{\pcIssueNeg}{0.18}
\newcommand{\pcLinked}{4{,}517}
\newcommand{\pcLinkedCostMed}{350{,}000}
\newcommand{\pcLinkedDays}{500}
\newcommand{\pcLinkedExit}{20.9}
\newcommand{\pcLinkedFy}{2011}
\newcommand{\pcLinkedIssued}{78.7}
\newcommand{\pcLinkedShare}{7.8}
\newcommand{\pcLmDaysMedA}{274}
\newcommand{\pcLmDaysMedB}{219}
\newcommand{\pcLmDraws}{5}
\newcommand{\pcLmExitA}{23.7}
\newcommand{\pcLmExitB}{34.5}
\newcommand{\pcLmLagMed}{320}
\newcommand{\pcLmLinkedAfter}{923}
\newcommand{\pcLmLinkedAlive}{3{,}354}
\newcommand{\pcLmLinkedWith}{3{,}587}
\newcommand{\pcLmMedGap}{55}
\newcommand{\pcLmMedSurvive}{28}
\newcommand{\pcLmSdDays}{0.04}
\newcommand{\pcLmSdDaysHi}{0.06}
\newcommand{\pcLmSdDaysLo}{0.03}
\newcommand{\pcLmSdExit}{-0.24}
\newcommand{\pcLmSdIssued}{0.29}
\newcommand{\pcLmSupport}{97.1}
\newcommand{\pcLmSurviveDays}{9}
\newcommand{\pcLmSurviveExit}{152}
\newcommand{\pcLmUnlinkedAlive}{14{,}926}
\newcommand{\pcMinCommissionPermits}{434}
\newcommand{\pcMinYearN}{50}
\newcommand{\pcMinisterialPermits}{31}
\newcommand{\pcNoFiled}{1{,}013}
\newcommand{\pcNoted}{17}
\newcommand{\pcOTC}{879{,}974}
\newcommand{\pcOTCShare}{76.6}
\newcommand{\pcOrdinanceUnits}{25}
\newcommand{\pcPermits}{1{,}149{,}352}
\newcommand{\pcPermitsRepeating}{116{,}986}
\newcommand{\pcPrimaryRows}{1{,}149{,}358}
\newcommand{\pcPropC}{2016}
\newcommand{\pcRetrieved}{2026-09-11}
\newcommand{\pcRows}{1{,}295{,}168}
\newcommand{\pcRowsRepeating}{145{,}816}
\newcommand{\pcRulesAgree}{94.6}
\newcommand{\pcRulesCompared}{1{,}149{,}352}
\newcommand{\pcSbFourTwoThree}{45}
\newcommand{\pcSbThirtyFive}{76}
\newcommand{\pcSdDaysCell}{0.43}
\newcommand{\pcSdDaysRaw}{0.71}
\newcommand{\pcSdDaysUnits}{0.44}
\newcommand{\pcSdExitCell}{-0.16}
\newcommand{\pcSdExitRaw}{-0.25}
\newcommand{\pcSdIssuedCell}{0.13}
\newcommand{\pcSdLogCostCell}{0.23}
\newcommand{\pcSdLogCostRaw}{0.54}
\newcommand{\pcSdLogCostUnits}{0.16}
\newcommand{\pcSdUnitsCell}{0.07}
\newcommand{\pcSdUnitsRaw}{0.12}
\newcommand{\pcSmoothYears}{3}
\newcommand{\pcSupportCells}{97.3}
\newcommand{\pcSupportUnits}{99.7}
\newcommand{\pcSystemCols}{4}
\newcommand{\pcTboth}{1{,}627}
\newcommand{\pcTone}{2{,}239}
\newcommand{\pcToneUniv}{1{,}890}
\newcommand{\pcTsfBands}{21, 100}
\newcommand{\pcTthreeOnly}{1{,}773}
\newcommand{\pcTthreePrecision}{0.72}
\newcommand{\pcTthreeRecall}{0.92}
\newcommand{\pcTtwo}{5{,}518}
\newcommand{\pcTtwoUniv}{4{,}042}
\newcommand{\pcUniverse}{57{,}883}
\newcommand{\pcUniverseShare}{5.0}
\newcommand{\pcUnlinkedCostMed}{150{,}000}
\newcommand{\pcUnlinkedDays}{191}
\newcommand{\pcUnlinkedExit}{31.7}
\newcommand{\pcUnlinkedFy}{2001}
\newcommand{\pcUnlinkedIssued}{81.5}
\newcommand{\pcYearBin}{5}

% GENERATED BY analyze_permit_content.py clocks --- do not edit by hand.
\newcommand{\ckAppHearingCvCu}{1.32}
\newcommand{\ckAppHearingCvDr}{1.27}
\newcommand{\ckAppHearingCvLpa}{1.12}
\newcommand{\ckAppHearingMedCu}{161}
\newcommand{\ckAppHearingMedDr}{301}
\newcommand{\ckAppHearingMedLpa}{603}
\newcommand{\ckAppHearingNCu}{3{,}141}
\newcommand{\ckAppHearingNDr}{2{,}311}
\newcommand{\ckAppHearingNLpa}{141}
\newcommand{\ckAppHearingNegCu}{6}
\newcommand{\ckAppHearingNegDr}{67}
\newcommand{\ckAppHearingNegLpa}{3}
\newcommand{\ckClockFrom}{1998}
\newcommand{\ckCvHiCu}{0.81}
\newcommand{\ckCvHiDr}{0.58}
\newcommand{\ckCvHiLpa}{---}
\newcommand{\ckCvHiMin}{0.98}
\newcommand{\ckCvLoCu}{0.67}
\newcommand{\ckCvLoDr}{0.71}
\newcommand{\ckCvLoLpa}{---}
\newcommand{\ckCvLoMin}{0.86}
\newcommand{\ckCvRhoCu}{0.10}
\newcommand{\ckCvRhoDr}{-0.40}
\newcommand{\ckCvRhoLpa}{---}
\newcommand{\ckCvRhoMin}{0.30}
\newcommand{\ckFiledIssuedCvCu}{1.19}
\newcommand{\ckFiledIssuedCvDr}{0.82}
\newcommand{\ckFiledIssuedCvLpa}{0.93}
\newcommand{\ckFiledIssuedCvMin}{1.08}
\newcommand{\ckFiledIssuedMedCu}{434}
\newcommand{\ckFiledIssuedMedDr}{631}
\newcommand{\ckFiledIssuedMedLpa}{722}
\newcommand{\ckFiledIssuedMedMin}{338}
\newcommand{\ckFiledIssuedNCu}{1{,}436}
\newcommand{\ckFiledIssuedNDr}{2{,}311}
\newcommand{\ckFiledIssuedNLpa}{115}
\newcommand{\ckFiledIssuedNMin}{28{,}498}
\newcommand{\ckFiledIssuedNegCu}{0}
\newcommand{\ckFiledIssuedNegDr}{1}
\newcommand{\ckFiledIssuedNegLpa}{0}
\newcommand{\ckFiledIssuedNegMin}{0}
\newcommand{\ckFinalIssuedCvCu}{1.88}
\newcommand{\ckFinalIssuedCvDr}{2.06}
\newcommand{\ckFinalIssuedCvLpa}{1.37}
\newcommand{\ckFinalIssuedMedCu}{408}
\newcommand{\ckFinalIssuedMedDr}{230}
\newcommand{\ckFinalIssuedMedLpa}{700}
\newcommand{\ckFinalIssuedNCu}{1{,}328}
\newcommand{\ckFinalIssuedNDr}{2{,}088}
\newcommand{\ckFinalIssuedNLpa}{94}
\newcommand{\ckFinalIssuedNegCu}{90}
\newcommand{\ckFinalIssuedNegDr}{70}
\newcommand{\ckFinalIssuedNegLpa}{4}
\newcommand{\ckFirstFinalPctCu}{73}
\newcommand{\ckFirstFinalPctDr}{68}
\newcommand{\ckFirstFinalPctLpa}{75}
\newcommand{\ckFirstFinalPctMin}{---}
\newcommand{\ckHearingFinalCvCu}{3.52}
\newcommand{\ckHearingFinalCvDr}{2.51}
\newcommand{\ckHearingFinalCvLpa}{3.02}
\newcommand{\ckHearingFinalMedCu}{0}
\newcommand{\ckHearingFinalMedDr}{0}
\newcommand{\ckHearingFinalMedLpa}{0}
\newcommand{\ckHearingFinalNCu}{3{,}205}
\newcommand{\ckHearingFinalNDr}{2{,}467}
\newcommand{\ckHearingFinalNLpa}{150}
\newcommand{\ckHearingFinalNegCu}{0}
\newcommand{\ckHearingFinalNegDr}{0}
\newcommand{\ckHearingFinalNegLpa}{0}
\newcommand{\ckHorizon}{5{,}500}
\newcommand{\ckIqrRhoCu}{---}
\newcommand{\ckIqrRhoDr}{---}
\newcommand{\ckIqrRhoLpa}{---}
\newcommand{\ckIqrRhoMin}{---}
\newcommand{\ckIssuedDoneCvCu}{1.02}
\newcommand{\ckIssuedDoneCvDr}{0.87}
\newcommand{\ckIssuedDoneCvLpa}{0.49}
\newcommand{\ckIssuedDoneCvMin}{0.98}
\newcommand{\ckIssuedDoneMedCu}{762}
\newcommand{\ckIssuedDoneMedDr}{759}
\newcommand{\ckIssuedDoneMedLpa}{1{,}307}
\newcommand{\ckIssuedDoneMedMin}{642}
\newcommand{\ckIssuedDoneNCu}{1{,}117}
\newcommand{\ckIssuedDoneNDr}{1{,}912}
\newcommand{\ckIssuedDoneNLpa}{90}
\newcommand{\ckIssuedDoneNMin}{22{,}233}
\newcommand{\ckIssuedDoneNegCu}{3}
\newcommand{\ckIssuedDoneNegDr}{2}
\newcommand{\ckIssuedDoneNegLpa}{0}
\newcommand{\ckIssuedDoneNegMin}{7}
\newcommand{\ckMature}{2013}
\newcommand{\ckMinCell}{20}
\newcommand{\ckPermitCvHiCu}{1.23}
\newcommand{\ckPermitCvHiDr}{0.88}
\newcommand{\ckPermitCvHiLpa}{---}
\newcommand{\ckPermitCvHiMin}{1.13}
\newcommand{\ckPermitCvLoCu}{1.12}
\newcommand{\ckPermitCvLoDr}{0.99}
\newcommand{\ckPermitCvLoLpa}{---}
\newcommand{\ckPermitCvLoMin}{1.18}
\newcommand{\ckPermitCvRhoCu}{0.10}
\newcommand{\ckPermitCvRhoDr}{-0.10}
\newcommand{\ckPermitCvRhoLpa}{---}
\newcommand{\ckPermitCvRhoMin}{-0.60}
\newcommand{\ckPermitPctCu}{45}
\newcommand{\ckPermitPctDr}{94}
\newcommand{\ckPermitPctLpa}{77}
\newcommand{\ckPermitPctMin}{100}
\newcommand{\ckRetrieved}{2026-09-11}
\newcommand{\ckTotCvCu}{0.79}
\newcommand{\ckTotCvDr}{0.67}
\newcommand{\ckTotCvLpa}{0.52}
\newcommand{\ckTotCvMin}{0.95}
\newcommand{\ckTotIqrCu}{---}
\newcommand{\ckTotIqrDr}{---}
\newcommand{\ckTotIqrLpa}{1.26}
\newcommand{\ckTotIqrMin}{---}
\newcommand{\ckTotMatNCu}{423}
\newcommand{\ckTotMatNDr}{955}
\newcommand{\ckTotMatNLpa}{40}
\newcommand{\ckTotMatNMin}{9{,}373}
\newcommand{\ckTotMedCu}{2{,}232}
\newcommand{\ckTotMedDr}{1{,}817}
\newcommand{\ckTotMedLpa}{3{,}394}
\newcommand{\ckTotMedMin}{1{,}334}
\newcommand{\ckTotPninetyCu}{---}
\newcommand{\ckTotPninetyDr}{---}
\newcommand{\ckTotPninetyLpa}{---}
\newcommand{\ckTotPninetyMin}{---}
\newcommand{\ckTotStenPctCu}{36}
\newcommand{\ckTotStenPctDr}{34}
\newcommand{\ckTotStenPctLpa}{48}
\newcommand{\ckTotStenPctMin}{34}
\newcommand{\ckTotalCvCu}{0.76}
\newcommand{\ckTotalCvDr}{0.64}
\newcommand{\ckTotalCvLpa}{0.50}
\newcommand{\ckTotalCvMin}{0.80}
\newcommand{\ckTotalMedCu}{2{,}232}
\newcommand{\ckTotalMedDr}{1{,}817}
\newcommand{\ckTotalMedLpa}{3{,}394}
\newcommand{\ckTotalMedMin}{1{,}334}
\newcommand{\ckTotalNCu}{1{,}440}
\newcommand{\ckTotalNDr}{2{,}311}
\newcommand{\ckTotalNLpa}{115}
\newcommand{\ckTotalNMin}{28{,}498}
\newcommand{\ckTotalNegCu}{0}
\newcommand{\ckTotalNegDr}{1}
\newcommand{\ckTotalNegLpa}{0}
\newcommand{\ckTotalNegMin}{7}
\newcommand{\ckUnitsCu}{3{,}205}
\newcommand{\ckUnitsDr}{2{,}467}
\newcommand{\ckUnitsLpa}{150}
\newcommand{\ckUnitsMin}{28{,}498}
\newcommand{\ckValueCover}{87.8}
\newcommand{\ckValueQ}{5}
\newcommand{\ckValuedLater}{14{,}593}

% GENERATED BY analyze_permit_content.py outcomes --- do not edit by hand.
\newcommand{\ocCifCu}{68}
\newcommand{\ocCifDr}{68}
\newcommand{\ocCifExitCu}{16}
\newcommand{\ocCifExitDr}{20}
\newcommand{\ocCifExitLpa}{7}
\newcommand{\ocCifExitMin}{28}
\newcommand{\ocCifLpa}{64}
\newcommand{\ocCifMin}{56}
\newcommand{\ocDoneCu}{79}
\newcommand{\ocDoneDr}{78}
\newcommand{\ocDoneHiCu}{77}
\newcommand{\ocDoneHiDr}{70}
\newcommand{\ocDoneHiMin}{60}
\newcommand{\ocDoneHp}{60}
\newcommand{\ocDoneLoCu}{81}
\newcommand{\ocDoneLoDr}{70}
\newcommand{\ocDoneLoMin}{65}
\newcommand{\ocDoneLpa}{94}
\newcommand{\ocDoneMin}{63}
\newcommand{\ocExitCu}{17}
\newcommand{\ocExitDr}{16}
\newcommand{\ocExitLpa}{3}
\newcommand{\ocExitMin}{28}
\newcommand{\ocFcdCover}{46}
\newcommand{\ocIssuedCu}{87}
\newcommand{\ocIssuedDr}{91}
\newcommand{\ocIssuedLpa}{97}
\newcommand{\ocIssuedMin}{75}
\newcommand{\ocMature}{2{,}136}
\newcommand{\ocMatureFrom}{2005}
\newcommand{\ocMatureTo}{2014}
\newcommand{\ocMedYearsCu}{3.7}
\newcommand{\ocMedYearsDr}{2.6}
\newcommand{\ocMedYearsLpa}{4.1}
\newcommand{\ocMedYearsMin}{2.8}
\newcommand{\ocNCu}{141}
\newcommand{\ocNDr}{171}
\newcommand{\ocNLpa}{33}
\newcommand{\ocNMin}{1{,}791}
\newcommand{\ocNeverCu}{32}
\newcommand{\ocNeverDr}{32}
\newcommand{\ocNeverLpa}{36}
\newcommand{\ocNeverMin}{44}
\newcommand{\ocOpenCu}{4}
\newcommand{\ocOpenDr}{6}
\newcommand{\ocOpenLpa}{3}
\newcommand{\ocOpenMin}{9}
\newcommand{\ocUnits}{10{,}253}
\newcommand{\ocUnitsApprCu}{6{,}583}
\newcommand{\ocUnitsApprDr}{386}
\newcommand{\ocUnitsApprLpa}{4{,}037}
\newcommand{\ocUnitsApprMin}{18{,}531}
\newcommand{\ocUnitsHpCu}{5{,}532}
\newcommand{\ocUnitsHpDr}{314}
\newcommand{\ocUnitsHpLpa}{3{,}952}
\newcommand{\ocUnitsHpMin}{8{,}715}
\newcommand{\ocUnitsShareCu}{84}
\newcommand{\ocUnitsShareDr}{81}
\newcommand{\ocUnitsShareLpa}{98}
\newcommand{\ocUnitsShareMin}{47}
\newcommand{\ocYears}{10}

% GENERATED BY build_exaction_panel.py report --- do not edit by hand.
\newcommand{\exAffordable}{188}
\newcommand{\exAlterations}{4{,}039}
\newcommand{\exAssumedGfa}{74}
\newcommand{\exBig}{843}
\newcommand{\exBigLinked}{380}
\newcommand{\exBunchEscPost}{942}
\newcommand{\exBunchEscPre}{880}
\newcommand{\exBunchEscRatio}{0.93}
\newcommand{\exBunchNoPost}{826}
\newcommand{\exBunchNoPre}{798}
\newcommand{\exBunchNoRatio}{0.97}
\newcommand{\exBunchTenPost}{70}
\newcommand{\exBunchTenPre}{83}
\newcommand{\exBunchTenRatio}{1.19}
\newcommand{\exBunchWeeks}{12}
\newcommand{\exClaims}{394}
\newcommand{\exClaimsUsed}{65}
\newcommand{\exClaimsVerified}{394}
\newcommand{\exCrossAgree}{53}
\newcommand{\exCrossAgreeTwo}{57}
\newcommand{\exCrossGf}{35}
\newcommand{\exCrossGfN}{17}
\newcommand{\exCrossLapsed}{46}
\newcommand{\exCrossLapsedN}{13}
\newcommand{\exCrossN}{121}
\newcommand{\exCrossNtwo}{107}
\newcommand{\exCrossRows}{260}
\newcommand{\exCrossSeventeen}{50}
\newcommand{\exCrossSeventeenN}{30}
\newcommand{\exCrossSix}{84}
\newcommand{\exCrossSixN}{19}
\newcommand{\exCrossThirteen}{52}
\newcommand{\exCrossThirteenN}{23}
\newcommand{\exCrossTwo}{69}
\newcommand{\exCrossTwoN}{13}
\newcommand{\exCutPost}{221}
\newcommand{\exCutPre}{214}
\newcommand{\exDbiApp}{26}
\newcommand{\exDensity}{49}
\newcommand{\exDensityBig}{45}
\newcommand{\exEffSix}{2006-09-09}
\newcommand{\exEffSixInferred}{2006-09-09}
\newcommand{\exElevenA}{3}
\newcommand{\exFiveSix}{4}
\newcommand{\exFourSix}{14}
\newcommand{\exGpu}{841}
\newcommand{\exGpuHi}{1{,}086}
\newcommand{\exGpuLo}{645}
\newcommand{\exGpuN}{1{,}284}
\newcommand{\exInclShare}{81}
\newcommand{\exLegalClaims}{112}
\newcommand{\exLinked}{966}
\newcommand{\exMedLag}{1.7}
\newcommand{\exNA}{131}
\newcommand{\exNB}{145}
\newcommand{\exNC}{99}
\newcommand{\exND}{25}
\newcommand{\exNineA}{18}
\newcommand{\exNineSix}{6}
\newcommand{\exPaySentenceAbsent}{2026}
\newcommand{\exPaySentenceFirst}{2011}
\newcommand{\exPaySentenceLast}{2025}
\newcommand{\exPipe}{94}
\newcommand{\exPriced}{400}
\newcommand{\exPricedLinked}{254}
\newcommand{\exProjects}{3{,}574}
\newcommand{\exPuHi}{\$78k}
\newcommand{\exPuLo}{\$50k}
\newcommand{\exPuMed}{\$61k}
\newcommand{\exPuMedA}{\$52k}
\newcommand{\exPuMedB}{\$71k}
\newcommand{\exPuMedC}{\$66k}
\newcommand{\exPuMedD}{\$64k}
\newcommand{\exPuMedLinked}{\$62k}
\newcommand{\exPuMedNot}{\$58k}
\newcommand{\exRegFirst}{2011}
\newcommand{\exRegLast}{2026}
\newcommand{\exRegisterClaims}{282}
\newcommand{\exRegisters}{15}
\newcommand{\exSbThreeThirty}{42}
\newcommand{\exShareHi}{28}
\newcommand{\exShareLo}{15}
\newcommand{\exShareMed}{20}
\newcommand{\exShareMedLinked}{21}
\newcommand{\exShareMedNot}{19}
\newcommand{\exStepTen}{\$420k}
\newcommand{\exStepTenPerUnit}{\$42k}
\newcommand{\exStepTsf}{\$209k}
\newcommand{\exStepTwentyFive}{\$525k}
\newcommand{\exStepTwentyFiveB}{\$26k}
\newcommand{\exStepYear}{2026}
\newcommand{\exTenA}{5}
\newcommand{\exTenSix}{2}
\newcommand{\exTenureGapMax}{\$20k}
\newcommand{\exTenureGapMed}{\$0k}
\newcommand{\exTfr}{124}
\newcommand{\exTwentyA}{4}
\newcommand{\exTwentyFiveA}{4}
\newcommand{\exTwentyFourA}{10}
\newcommand{\exTwentyOneA}{1}
\newcommand{\exUmu}{49}
\newcommand{\exVersions}{8}
\newcommand{\exWindow}{59}
\makeatletter
\expandafter\def\csname cl@inc2002_apply\endcsname{2001-06-18}
\expandafter\def\csname cl@inc2002_threshold\endcsname{10}
\expandafter\def\csname cl@inc2002_onsite_10\endcsname{10}
\expandafter\def\csname cl@inc2002_onsite_12\endcsname{12}
\expandafter\def\csname cl@inc2002_offsite_15\endcsname{15}
\expandafter\def\csname cl@inc2002_offsite_17\endcsname{17}
\expandafter\def\csname cl@inc2002_fee_basis\endcsname{off-site unit count x affordability gap}
\expandafter\def\csname cl@inc2002_transition\endcsname{5}
\expandafter\def\csname cl@inc2002_effective\endcsname{2002-04-05}
\expandafter\def\csname cl@inc2006_5to9\endcsname{2006-07-18}
\expandafter\def\csname cl@inc2006_onsite_15\endcsname{15}
\expandafter\def\csname cl@inc2006_offsite_20\endcsname{20}
\expandafter\def\csname cl@inc2006_vesting\endcsname{first site/building permit}
\expandafter\def\csname cl@inc2006_passed\endcsname{2006-08-01}
\expandafter\def\csname cl@inc2006_sep9\endcsname{2006-09-09}
\expandafter\def\csname cl@inc2006_enacted\endcsname{2006-08-10}
\expandafter\def\csname cl@inc2013_table_fee20\endcsname{20}
\expandafter\def\csname cl@inc2013_table_onsite12\endcsname{12}
\expandafter\def\csname cl@inc2013_table_offsite20\endcsname{20}
\expandafter\def\csname cl@inc2013_table_july18\endcsname{2006-07-18}
\expandafter\def\csname cl@inc2013_threshold10\endcsname{2013-01-15}
\expandafter\def\csname cl@inc2013_table_disclaimer\endcsname{yes}
\expandafter\def\csname cl@inc2013_effective\endcsname{2013-05-10}
\expandafter\def\csname cl@inc2016_fee20\endcsname{20}
\expandafter\def\csname cl@inc2016_fee33\endcsname{33}
\expandafter\def\csname cl@inc2016_onsite\endcsname{12 / 25}
\expandafter\def\csname cl@inc2016_fee_timing\endcsname{first construction document}
\expandafter\def\csname cl@propc2016\endcsname{2016-06-07}
\expandafter\def\csname cl@propc2016_contingent\endcsname{Charter 16.110}
\expandafter\def\csname cl@propc2012\endcsname{2012-11}
\expandafter\def\csname cl@ord_effective_30days\endcsname{30 days}
\expandafter\def\csname cl@inc2016_effective\endcsname{2016-06-12}
\expandafter\def\csname cl@inc2017_effective\endcsname{2017-08-26}
\expandafter\def\csname cl@inc2017_cutoff\endcsname{2016-01-12}
\expandafter\def\csname cl@inc2017_after2016\endcsname{2016-01-12}
\expandafter\def\csname cl@cod_4153_pre2016\endcsname{2016-01-12}
\expandafter\def\csname cl@cod_4153_25plus\endcsname{2013-01-01}
\expandafter\def\csname cl@cod_4153_on13\endcsname{13}
\expandafter\def\csname cl@cod_4153_on135\endcsname{13.5}
\expandafter\def\csname cl@cod_4153_on145\endcsname{14.5}
\expandafter\def\csname cl@cod_4153_fee25\endcsname{25}
\expandafter\def\csname cl@cod_4153_fee275\endcsname{27.5}
\expandafter\def\csname cl@cod_4153_fee30\endcsname{30}
\expandafter\def\csname cl@cod_4153_deadline\endcsname{2018-12-07}
\expandafter\def\csname cl@cod_4153_exempt\endcsname{2016-01-12}
\expandafter\def\csname cl@cod_4153_threshold\endcsname{10}
\expandafter\def\csname cl@cod_4155_fee20\endcsname{20}
\expandafter\def\csname cl@cod_4155_fee33\endcsname{33}
\expandafter\def\csname cl@cod_4155_fee30\endcsname{30}
\expandafter\def\csname cl@cod_4155_vesting\endcsname{complete EEA date}
\expandafter\def\csname cl@cod_4155_30months\endcsname{30 months}
\expandafter\def\csname cl@cod_4155_timing\endcsname{402(d)}
\expandafter\def\csname cl@cod_4155_dbl\endcsname{yes}
\expandafter\def\csname cl@cod_4156_on12\endcsname{12}
\expandafter\def\csname cl@cod_4156_on20\endcsname{20}
\expandafter\def\csname cl@cod_4156_on18\endcsname{18}
\expandafter\def\csname cl@cod_4156_esc_small\endcsname{0.5}
\expandafter\def\csname cl@cod_4156_esc_large1\endcsname{1.0}
\expandafter\def\csname cl@cod_4156_esc_large2\endcsname{0.5}
\expandafter\def\csname cl@cod_4156_caps\endcsname{26 / 24}
\expandafter\def\csname cl@cod_4156_vesting\endcsname{complete EEA date}
\expandafter\def\csname cl@cod_4156b_same\endcsname{20 / 18}
\expandafter\def\csname cl@inc2023_effective\endcsname{2023-11-12}
\expandafter\def\csname cl@a415_pipeline\endcsname{2023-11-01}
\expandafter\def\csname cl@a415_final_approval\endcsname{first development application}
\expandafter\def\csname cl@a415_request\endcsname{2026-11-01}
\expandafter\def\csname cl@a415_fee164\endcsname{16.4}
\expandafter\def\csname cl@a415_onsite12_small\endcsname{12}
\expandafter\def\csname cl@a415_onsite12_large\endcsname{12}
\expandafter\def\csname cl@a415_offsite164\endcsname{16.4}
\expandafter\def\csname cl@b415_window\endcsname{2023-11-01 to 2026-11-01}
\expandafter\def\csname cl@b415_fee205\endcsname{20.5}
\expandafter\def\csname cl@b415_onsite15\endcsname{15}
\expandafter\def\csname cl@area_specific_rates\endcsname{area-specific}
\expandafter\def\csname cl@b415_sunset\endcsname{2026-11-01}
\expandafter\def\csname cl@o201_operative\endcsname{2026-11-21}
\expandafter\def\csname cl@o201_2026_fee27\endcsname{27}
\expandafter\def\csname cl@o201_2026_fee245\endcsname{24.5}
\expandafter\def\csname cl@o201_2026_onsite15\endcsname{15}
\expandafter\def\csname cl@o201_2026_fee20\endcsname{20}
\expandafter\def\csname cl@o201_2026_on20o\endcsname{20}
\expandafter\def\csname cl@o201_2026_on18r\endcsname{18}
\expandafter\def\csname cl@o201_2026_off27o\endcsname{27}
\expandafter\def\csname cl@o201_2026_esc\endcsname{2028-01-01}
\expandafter\def\csname cl@fee_402d_timing\endcsname{first construction document}
\expandafter\def\csname cl@fee_409_index\endcsname{January 1}
\expandafter\def\csname cl@fee_rate_at_payment\endcsname{payment date}
\expandafter\def\csname cl@fee_193_approval\endcsname{final approval}
\expandafter\def\csname cl@fee_193_effective\endcsname{2023-10-16}
\expandafter\def\csname cl@fee_tfr33\endcsname{33}
\expandafter\def\csname cl@fee_tfr_tsf\endcsname{411A}
\expandafter\def\csname cl@tsf_threshold\endcsname{21}
\expandafter\def\csname cl@tsf_base\endcsname{7.74 / 8.74}
\expandafter\def\csname cl@tsf_created\endcsname{2015-12-25}
\expandafter\def\csname cl@cc_threshold\endcsname{1}
\expandafter\def\csname cl@cc_created\endcsname{2016-02-18}
\expandafter\def\csname cl@cc_base\endcsname{1.83 / 0.91}
\expandafter\def\csname cl@jhl_nonres\endcsname{non-residential}
\expandafter\def\csname cl@area_en_start\endcsname{2008-12-19}
\expandafter\def\csname cl@area_mo421_start\endcsname{2008-04-03}
\expandafter\def\csname cl@area_mo416_start\endcsname{2008-05-30}
\expandafter\def\csname cl@area_rh_start\endcsname{2005-08-19}
\expandafter\def\csname cl@area_soma_start\endcsname{2005-08-19}
\expandafter\def\csname cl@area_vv_start\endcsname{2005-11-18}
\expandafter\def\csname cl@area_bp_start\endcsname{2009-04-17}
\expandafter\def\csname cl@area_tc_start\endcsname{2012-09-07}
\expandafter\def\csname cl@area_cs_start\endcsname{2019-01-12}
\expandafter\def\csname cl@sb330_freeze\endcsname{preliminary application}
\expandafter\def\csname cl@sb330_index_exception\endcsname{cost index}
\expandafter\def\csname cl@sb330_chapter\endcsname{Stats. 2019, ch. 654}
\expandafter\def\csname cl@sb330_approved\endcsname{2019-10-09}
\expandafter\def\csname cl@gov_standards_fees\endcsname{yes}
\makeatother
\newcommand{\cl}[1]{\csname cl@#1\endcsname}

% GENERATED BY audit_state_laws.py report --- do not edit by hand.
\newcommand{\rtCa}{29}
\newcommand{\rtCasesPost}{935}
\newcommand{\rtCasesPre}{964}
\newcommand{\rtCatFive}{17}
\newcommand{\rtCatFour}{3}
\newcommand{\rtCatOne}{9}
\newcommand{\rtCatThree}{6}
\newcommand{\rtCatTwo}{10}
\newcommand{\rtClaims}{156}
\newcommand{\rtClaimsOk}{156}
\newcommand{\rtEffAbEightEightOne}{2020-01-01}
\newcommand{\rtEffAbFiveEightSeven}{2020-01-01}
\newcommand{\rtEffAbOneFiveOneFive}{2018-01-01}
\newcommand{\rtEffAbOneSevenSixThree}{2020-01-01}
\newcommand{\rtEffAbOneSixThreeThree}{2024-01-01}
\newcommand{\rtEffAbOneThreeZero}{2025-06-30}
\newcommand{\rtEffAbSixEight}{2020-01-01}
\newcommand{\rtEffAbSixSevenEight}{2018-01-01}
\newcommand{\rtEffAbTwoThreeFourFive}{2021-01-01}
\newcommand{\rtEffAbTwoTwoFourThree}{2025-01-01}
\newcommand{\rtEffAbTwoTwoTwoOne}{2023-01-01}
\newcommand{\rtEffAbTwoZeroNineSeven}{2023-01-01}
\newcommand{\rtEffAbTwoZeroOneOne}{2023-01-01}
\newcommand{\rtEffCuappeal}{2022-10-17}
\newcommand{\rtEffDbl}{---}
\newcommand{\rtEffFeedefer}{2025-11-24}
\newcommand{\rtEffFhosud}{2023-10-16}
\newcommand{\rtEffFzp}{2026-01-12}
\newcommand{\rtEffHaa}{---}
\newcommand{\rtEffHeTwoZeroTwoTwo}{2023-03-03}
\newcommand{\rtEffHomesf}{2017-07-13}
\newcommand{\rtEffHpTwoZeroTwoFour}{2024-04-22}
\newcommand{\rtEffHpTwoZeroTwoThree}{2024-01-14}
\newcommand{\rtEffIncwaiver}{2026-01-23}
\newcommand{\rtEffInferred}{18}
\newcommand{\rtEffJulyTwoZeroTwoSix}{---}
\newcommand{\rtEffMindens}{2025-01-19}
\newcommand{\rtEffNotes}{5}
\newcommand{\rtEffOrdOneFiveEight}{2017-08-26}
\newcommand{\rtEffOrdOneNineThree}{2023-10-16}
\newcommand{\rtEffOrdTwoZeroOne}{2023-11-12}
\newcommand{\rtEffPropc}{2016-06-12}
\newcommand{\rtEffPsa}{---}
\newcommand{\rtEffRhna}{---}
\newcommand{\rtEffSbEight}{2022-01-01}
\newcommand{\rtEffSbEightNineSeven}{2023-01-01}
\newcommand{\rtEffSbFour}{2024-01-01}
\newcommand{\rtEffSbFourTwoThree}{2024-01-01}
\newcommand{\rtEffSbNine}{2022-01-01}
\newcommand{\rtEffSbOneSixSeven}{2018-01-01}
\newcommand{\rtEffSbOneThreeOne}{2025-06-30}
\newcommand{\rtEffSbOneZero}{2022-01-01}
\newcommand{\rtEffSbSevenNine}{2026-01-01}
\newcommand{\rtEffSbSix}{2023-01-01}
\newcommand{\rtEffSbThreeFive}{2018-01-01}
\newcommand{\rtEffSbThreeThreeZero}{2020-01-01}
\newcommand{\rtEffSitepermit}{2023-08-28}
\newcommand{\rtEnactedAbEightEightOne}{2019-10-09}
\newcommand{\rtEnactedAbFiveEightSeven}{2019-10-09}
\newcommand{\rtEnactedAbOneFiveOneFive}{2017-09-29}
\newcommand{\rtEnactedAbOneSevenSixThree}{2019-10-09}
\newcommand{\rtEnactedAbOneSixThreeThree}{2023-10-11}
\newcommand{\rtEnactedAbOneThreeZero}{2025-06-30}
\newcommand{\rtEnactedAbSixEight}{2019-10-09}
\newcommand{\rtEnactedAbSixSevenEight}{2017-09-29}
\newcommand{\rtEnactedAbTwoThreeFourFive}{2020-09-28}
\newcommand{\rtEnactedAbTwoTwoFourThree}{2024-09-19}
\newcommand{\rtEnactedAbTwoTwoTwoOne}{2022-09-28}
\newcommand{\rtEnactedAbTwoZeroNineSeven}{2022-09-22}
\newcommand{\rtEnactedAbTwoZeroOneOne}{2022-09-28}
\newcommand{\rtEnactedCuappeal}{2022-09-06}
\newcommand{\rtEnactedDbl}{---}
\newcommand{\rtEnactedFeedefer}{2025-10-21}
\newcommand{\rtEnactedFhosud}{2023-09-05}
\newcommand{\rtEnactedFzp}{2025-12-09}
\newcommand{\rtEnactedHaa}{---}
\newcommand{\rtEnactedHeTwoZeroTwoTwo}{2023-01-31}
\newcommand{\rtEnactedHomesf}{2017-06-06}
\newcommand{\rtEnactedHpTwoZeroTwoFour}{2024-03-12}
\newcommand{\rtEnactedHpTwoZeroTwoThree}{2023-12-12}
\newcommand{\rtEnactedIncwaiver}{2025-12-16}
\newcommand{\rtEnactedJulyTwoZeroTwoSix}{---}
\newcommand{\rtEnactedMindens}{2024-12-17}
\newcommand{\rtEnactedOrdOneFiveEight}{2017-07-18}
\newcommand{\rtEnactedOrdOneNineThree}{2023-09-05}
\newcommand{\rtEnactedOrdTwoZeroOne}{2023-10-03}
\newcommand{\rtEnactedPropc}{2016-05-03}
\newcommand{\rtEnactedPsa}{---}
\newcommand{\rtEnactedRhna}{---}
\newcommand{\rtEnactedSbEight}{2021-09-16}
\newcommand{\rtEnactedSbEightNineSeven}{2022-09-28}
\newcommand{\rtEnactedSbFour}{2023-10-11}
\newcommand{\rtEnactedSbFourTwoThree}{2023-10-11}
\newcommand{\rtEnactedSbNine}{2021-09-16}
\newcommand{\rtEnactedSbOneSixSeven}{2017-09-29}
\newcommand{\rtEnactedSbOneThreeOne}{2025-06-30}
\newcommand{\rtEnactedSbOneZero}{2021-09-16}
\newcommand{\rtEnactedSbSevenNine}{2025-10-10}
\newcommand{\rtEnactedSbSix}{2022-09-28}
\newcommand{\rtEnactedSbThreeFive}{2017-09-29}
\newcommand{\rtEnactedSbThreeThreeZero}{2019-10-09}
\newcommand{\rtEnactedSitepermit}{2023-07-25}
\newcommand{\rtFirstQ}{1998Q1}
\newcommand{\rtHearWinPost}{2020--2025}
\newcommand{\rtHearWinPre}{2015--2019}
\newcommand{\rtItemsLast}{2026-03-26}
\newcommand{\rtItemsPost}{91}
\newcommand{\rtItemsPre}{111}
\newcommand{\rtLastQ}{2026Q1}
\newcommand{\rtLaws}{45}
\newcommand{\rtLinkWinEarly}{2005--2009}
\newcommand{\rtLinkWinLate}{2020--2025}
\newcommand{\rtLinkedEarly}{13.1}
\newcommand{\rtLinkedLate}{14.3}
\newcommand{\rtNotMeasurable}{28}
\newcommand{\rtOptIn}{19}
\newcommand{\rtOverFiveNPost}{22}
\newcommand{\rtOverFiveNPre}{18}
\newcommand{\rtOverFivePost}{2.4}
\newcommand{\rtOverFivePre}{1.9}
\newcommand{\rtPartUsable}{11}
\newcommand{\rtPartial}{0}
\newcommand{\rtPopAbTwentyEleven}{AB 2011 in text: 5 a year, 2021--2025}
\newcommand{\rtPopAdu}{DBI ADU permits: 182 a year, 2021--2025}
\newcommand{\rtPopDbl}{density bonus in text: 58 a year, 2021--2025}
\newcommand{\rtPopHomesf}{HOME-SF in text: 7 a year, 2021--2025}
\newcommand{\rtPopSbFourTwoThree}{SB 423 in text: 7 a year, 2021--2025}
\newcommand{\rtPopSbNine}{SB 9 or lot split in text: 15 a year, 2021--2025}
\newcommand{\rtPopSbThirtyFive}{SB 35 flag + text: 13 a year, 2021--2025}
\newcommand{\rtPopSbThreeThirty}{SB 330 flag (PRJ): 27 a year, 2021--2025}
\newcommand{\rtResCuPost}{22.7}
\newcommand{\rtResCuPre}{29.4}
\newcommand{\rtResDrPost}{13.2}
\newcommand{\rtResDrPre}{19.4}
\newcommand{\rtSf}{16}
\newcommand{\rtSfLastEff}{2026-08-23}
\newcommand{\rtSfLastOrd}{0133-26}
\newcommand{\rtSfMatch}{0}
\newcommand{\rtSfTwentySix}{131}
\newcommand{\rtShortlist}{SB 330 (Housing Crisis Act); Proposition C and Ord. 76-16; Ord. 158-17 (inclusionary); Ord. 201-23 (fee reductions); Housing Production (Ord. 248-23); Housing Production (Ord. 53-24)}
\newcommand{\rtUnverified}{1}
\newcommand{\rtUsable}{6}
\newcommand{\rtVerified}{44}
\newcommand{\rtWinPost}{2024--2026}
\newcommand{\rtWinPre}{2021--2022}

\newcommand{\chtable}[1]{\FloatBarrier\def\chtab{#1}% GENERATED BY build_chronicle.py --- do not edit by hand.
\ifnum\pdfstrcmp{\chtab}{timeline}=0
{\scriptsize\begin{longtable}{L{1.55cm}L{1.85cm}L{2.8cm}L{3.7cm}L{4.2cm}}
\caption{Every memo, in order of completion: the question it was asked, its answer in a sentence, and what later work revised or overturned in it.}\label{tab:chtimeline}\\\toprule
Date & Memo & Question & Answer & Revised or overturned by\\\midrule\endfirsthead
\toprule Date & Memo & Question & Answer & Revised or overturned by\\\midrule\endhead
\chDateMeetingLevelInfo & Meeting-level information (\S\ref{sec:chmeetinglevelinfo}) & Can the hearing date and the sitting's attributes be derived from the minutes? & Dates are right for every gold item; the attribute fields score in the mid-90s out of sample, worst in 1998--2001. & 2026-09-04: the 2018 scraping gap refilled. 2026-09-10: the round-4 figure (held-out meetings) and the all-meetings figure (\chMeetAcc\%, unadjudicated) are different measurements; a name-reducer defect had split six commissioners into two each.\\
\chDateExtractionMethodComparison & Item-level extraction (\S\ref{sec:chextractionmethodcomparison}) & Rules or a language model for the item fields? & Haiku with retrieved examples, \chHaikuTest\% against the regex's \chRegexTest\% on the frozen test half. & 2026-09-10: the regex baseline had been understated by its own scoring; few-shot retrieval had let an item answer itself, which touched every ``all items'' figure but not the test half.\\
\chDateDiscretionaryReviewPatterns & The corpus memo (\S\ref{sec:chdiscretionaryreviewpatterns}) & What does the item-level dataset contain? & Denial has nearly vanished while conditioning tripled; delay's upper tail lengthened; DR's share of the docket collapsed. & 2026-09-10: the conditioning rate was recomputed after a defect fix. 2026-09-11: the case universe fixed at \corpusUniverse; its distinct-parcel count is superseded by the acquisition memo's join.\\
\chDatePermitLinkage & Permit linkage (\S\ref{sec:chpermitlinkage}) & How far can a decision be followed into DBI's permits? & The printed permit number matches on \chPermitMatch\% of \chPermitsDistinct\ numbers; the parcel fallback runs at \chParcelPrecision\ precision; the Commission's conditions are not in DBI. & 2026-09-11: the fallback's precision was measured on discretionary reviews only, so the permit content memo keeps that tier (T3) apart; the Planning records' permit field is a second route.\\
\chDateConditionsOfApproval & Conditions of approval (\S\ref{sec:chconditionsofapproval}) & Where do the conditions of approval live, and how much is obtainable? & In the Commission packet's draft motion, addressable by case number from 2011; templated, and a thirteen-category scheme spans them. & 2026-09-08 and 2026-09-11: ``the two routes never cover the same item'' withdrawn (the records' permit field reaches \condBridgeCUPct\% of conditional uses); ``no amount of work recovers condition content before 2010'' withdrawn --- it was a claim about two URL patterns, and the census found text for \condReachEarlyHtml\% and \condReachLateHtml\% of the pre-2010 eras: little, not none.\\
\chDatePredictingDelay & Predicting delay (\S\ref{sec:chpredictingdelay}) & How much of the hearing-to-hearing time is forecastable? & The level somewhat, the tail barely; the residual dispersion is the first proxy for the cost of uncertainty. & 2026-09-11: the case universe note. Its clock starts at the first hearing; the acquisition memo's filing clocks and the clocks memo's full distributions supersede it as the duration measure.\\
\chDateDataAcquisition & Data acquisition (\S\ref{sec:chdataacquisition}) & Where are the denominator, the clock, the scale, the price and the fees? & A defined parcel-year risk set, a filing clock, a price on every parcel-year, and fees that vary across sections rather than years. & 2026-09-08: two of its own conclusions reversed (the lettered-lot key; the two-digit year). 2026-09-11: on the DR track the developer's clock starts at the DBI filing, not the Planning record (from the permit content memo). 2026-09-11: the feasible envelope added.\\
\chDateDatasfPlanningCatalogue & DataSF catalogue (\S\ref{sec:chdatasfplanningcatalogue}) & What does DataSF hold under ``planning'' that the pipeline does not use? & No dataset of Commission actions exists; review-time metrics and DBI's station routing are the unused delay data. & None yet.\\
\chDateConditionsContent & Conditions census (\S\ref{sec:chconditionscontent}) & What do the conditions say, across every case and motion? & Condition text for \ccCondWithTextPct\% of conditioned items; a template with a thin project-specific layer; the entitlement's term is a \ccValidityModal-month template. & 2026-09-11 (this brief): the parser's third round; the document once read as the earliest motion in the universe was two motions bound together, and every span resting on it moved; the numbers inside the conditions extracted.\\
\chDatePermitsContent & Permit content (\S\ref{sec:chpermitscontent}) & Are the projects that reach the Commission different from similar ones? & Bigger and slower, but most of the time difference is immortal time; the lower abandonment survives. & 2026-09-11 (this brief): the landmark comparison quantified how little of the time gap survives; cost placeholders dropped; the first-construction-document date dropped as an outcome.\\
\chDateClocks & Clocks (\S\ref{sec:chclocks}) & What are the full distributions of every duration, by track and parcel value? & Long and dispersed on every track; about a third never record completion; dispersion does not fall consistently with value. & Extended the same day with the outcome chain.\\
\chDateExactions & Exactions (\S\ref{sec:chexactions}) & What inclusionary requirement and impact fees vested for each project? & A sourced version chain and vesting rules; a median \exPuMed\ per unit for new buildings of ten or more units. & None yet; its open questions are stated in it.\\
\chDateRegulatoryTimeline & Regulatory timeline (\S\ref{sec:chregulatorytimeline}) & Which state laws and ordinances generate variation this data can measure? & \rtUsable\ of \rtLaws; the opt-in pathways are elections, not assignments. & None yet.\\
\bottomrule\end{longtable}}
\fi
\ifnum\pdfstrcmp{\chtab}{inventory}=0
{\scriptsize\begin{longtable}{L{2.9cm}L{3.0cm}L{2.0cm}L{4.4cm}L{2.0cm}}
\caption{The data inventory: every table and panel, its unit and size, its years, its known defects, and the memo that documents it.}\label{tab:chinventory}\\\toprule
Object & Unit (rows) & Years & Known defects & Memo\\\midrule\endfirsthead
\toprule Object & Unit (rows) & Years & Known defects & Memo\\\midrule\endhead
Minutes corpus and item table & item (\chItems) & \chItemFirst--\chItemLast & fields extracted by model at a measured error; modern era thin in the gold set & corpus; extraction\\
Meeting table & meeting (\chMeetings) & \chMeetFirst--\chMeetLast & attributes weakest in 1998--2001 (\chMeetEarly\% against \chMeetLate\% from 2015) & meeting level\\
Case universe & case (\corpusUniverse) & \chItemFirst--\chItemLast & case numbers as printed, two-digit years expanded & corpus; acquisition\\
Permit matches & printed permit number (\chPermitsDistinct) & \chItemFirst--\chItemLast & the printed route reaches DR, almost never CU & permit linkage\\
DBI permits, every column & permit (\pcPermits) & \pcFiledFirst--\pcFiledLast & cost placeholders; first-construction date sparse & permit content\\
Linkage tiers T1/T2/T3 & item--permit link & \chItemFirst--\chItemLast & T3 precision measured on DR only & permit content\\
Planning records & record & \chRecFirst--\chRecLast\ (sparse early) & record open date is not acceptance & acquisition\\
Parcel-year zoning panel & parcel-year (\acqPanelRows) & \acqPanelFirst--\acqPanelLast & zoning after \acqSnapshotFirst\ carried forward & acquisition\\
Risk set & parcel-year (\acqRiskParcelYears) & \acqPanelFirst--\acqPanelLast & condominium rule drops \acqRiskCondoLost\ hearings & acquisition\\
Price panel & parcel-year & \chPriceFirst--\chPriceLast & imputed where no sale, flagged & acquisition\\
Fee registers & rate by year & \exRegFirst--\exRegLast\ (no 2017) & pre-2011 missing & acquisition; exactions\\
Conditions census and parse & condition row (\chRowsConditions) & \chCondFirst--\chCondLast\ (thin early) & pre-2010 documents largely not online & conditions census\\
Conditions numbers & condition (\chRowsNumeric) & as the census & numeric fields only where stated & conditions census\\
Clocks panel & case or permit (\chRowsClocks) & \chClockFirst--\chClockLast & right-censoring; non-completion truncation & clocks\\
Outcome chain & housing project (\chRowsOutcomes) & \chClockFirst--\chClockLast & units delivered are a lower bound & clocks\\
Exaction panel & new-construction project (\chRowsExactions) & \chExFirst--\chExLast & priced from \exRegFirst; floor area mostly assumed; tenure mostly unknown & exactions\\
Law inventory and series & law (\chLaws) & enacted \chLawFirst--\chLawLast & \rtEffInferred\ effective dates inferred & regulatory timeline\\
Feasible envelope & parcel-year (\chRowsEnvelope) & \acqPanelFirst--\acqPanelLast & rules held at one reading; some districts unread & acquisition\\
Records request package & item (\chRrWith\ + \chRrWithout) & \chRrFirst--\chRrLast & draft; not sent & records request\\
\bottomrule\end{longtable}}
\fi
\ifnum\pdfstrcmp{\chtab}{ident}=0
\begin{table}[htbp]\centering\scriptsize
\caption{Which data object identifies which model object.}\label{tab:chident}
\begin{tabular}{L{3.5cm}L{4.2cm}L{3.6cm}L{2.9cm}}\toprule
Model primitive & Identifying moment & Dataset & Exists?\\\midrule
$\tau$: fees and inclusionary obligation & dollar obligation per project, by vesting date & exaction panel & yes (priced from 2011)\\
$\tau$: the template conditions & count and kind of conditions per motion & conditions census & yes (thin before 2010)\\
$\sigma$: delay dispersion & CV and Kaplan--Meier quantiles of each clock, by track and value & clocks panel & yes (cross-sectional, not ex ante)\\
$\sigma$: outcome risk & probability of denial, conditioning, abandonment & item table; outcome chain & yes\\
Arrival of discretion & share of projects that face a hearing, by parcel & risk set + linkage tiers & partly (T3 unvalidated outside DR)\\
Filing threshold & filing hazard by parcel value & risk set, price panel, planning records & yes\\
Parcel value & price per parcel-year & price panel & yes (imputed where no sale)\\
Capacity (the risk set's ``could have filed'') & maximum units per parcel-year & feasible envelope & yes (assumptions stated)\\
Option term & validity period and extensions & conditions numbers & yes\\
Separating $\tau$ from $\sigma$ & variation in $\tau$ with $\sigma$ fixed, and the reverse & inclusionary changes; SB 330; Housing Production ordinances & described, not estimated\\
Deterred projects & non-filing on the risk set & risk set & defined; selection not solved\\
\bottomrule\end{tabular}\end{table}
\fi
\ifnum\pdfstrcmp{\chtab}{register}=0
\begin{table}[htbp]\centering\scriptsize
\caption{The open register: what is unresolved, who or what resolves it, and whether it blocks a first presentable draft.}\label{tab:chregister}
\begin{tabular}{L{9.8cm}L{2.7cm}l}\toprule Open item & Needs & Blocks draft\\\midrule
Deterred projects: the risk set models non-filing, but who would have filed is unobserved & modelling decision & yes\\
The unconstrained-envelope counterfactual: an envelope exists; its relaxation is unspecified & modelling decision & no\\
Identification of $\tau$ separately from $\sigma$: the candidate variation is described, not estimated & modelling decision & yes\\
$\sigma$ is a parcel-level primitive estimated from a hazard --- an assumption, not a finding & modelling decision & yes\\
The pre-2010 motion residual & records request & no\\
Recorded Notices of Special Restrictions: document type, bulk access, pre-1990 terms & records request & no\\
Whether T3 linkage can be validated outside discretionary review & data & no\\
Area-specific inclusionary rates, TSF grandfathering, existing-use credits, the 2017 register & data & no\\
How the 2016 grandfathering deadline was applied (conditions contradict the codified text) & data & no\\
Effective dates of \rtEffInferred\ statutes inferred, not read & literature check & no\\
The Priority Equity Geographies polygon, for the Housing Production ordinances' treated set & data & no\\
The meeting-level gold set is unadjudicated for its all-meetings figure & data & no\\
The 25-unit ordinance of July 2026 named in the brief & literature check & no\\
Review-time metrics and DBI station routing (DataSF), unused & data & no\\
\bottomrule\end{tabular}\end{table}
\fi
\ifnum\pdfstrcmp{\chtab}{supersessions}=0
{\scriptsize\begin{longtable}{L{1.55cm}L{2.2cm}L{4.8cm}L{5.9cm}}
\caption{Every supersession found: what a memo claimed or measured, and what replaced it. Old and new figures are read from the logs and the brief that record them.}\label{tab:chsupersessions}\\\toprule
Date & Memo & Was & Now\\\midrule\endfirsthead\toprule Date & Memo & Was & Now\\\midrule\endhead
2026-09-04 & Meeting level & 2018 was a corpus hole of 2 documents; series crossed it & refilled from the S3 packet prefix to 42 documents\\
2026-09-07 & Extraction & accuracy against the hand labels as the measure & 52\% of 187 judged disagreements were label errors; unadjudicated accuracy is a lower bound\\
2026-09-07 & Extraction & gold records complete & a scalar/list mismatch had silently dropped lot numbers on 122 of 232 gold records; repaired\\
2026-09-08 & Permit linkage & the permit route collapsed after 2019 & a change in how the minutes print the number; reading it recovers 711 items and 483 permits\\
2026-09-08 & Conditions of approval & discretionary review carries no conditions section & 1 of 41 discretionary reviews do\\
2026-09-08 & Conditions of approval & the two linkage routes never cover the same item & the Planning records' permit field reaches \condBridgeCUPct\% of conditional uses\\
2026-09-08 & Data acquisition & parcel join below the acceptance criterion, at 89.5\% & 97.7\% once lettered lots and multi-block items are keyed correctly\\
2026-09-08 & Data acquisition & pre-2002 records do not go back that far under any key & expanding the two-digit year recovers 527 of 530 pre-2002 cases\\
2026-09-08 & Data acquisition & a first price surface & it carried a +0.90 log bias, caught by held-out scoring and rebuilt\\
2026-09-10 & Extraction & the regex baseline at 70.7\% & 75.2\% once scored through the same storage layer as the model\\
2026-09-10 & Meeting level & 146 distinct roll-call names & 78 once a mojibake token stopped splitting commissioners in two\\
2026-09-10 & Permit linkage and the analyses & the parcel-key fix applied & the lettered-lot bug was still live in four scripts; the fix grows the fallback's validation set from 2{,}909 to 3{,}274\\
2026-09-11 & Conditions of approval & no condition content before 2010 & \condReachEarlyHtml\% and \condReachLateHtml\% of the pre-2010 eras' conditioned items have text\\
2026-09-11 & The corpus memo & 8{,}295 distinct parcels & 8{,}980 from the acquisition memo's join (not yet regenerated)\\
2026-09-11 & Conditions census & the earliest motion read for a case in the universe was heard in \ccRprevEarliestInUniverse & two motions bound together; the earliest such item now dates from \ccRcurEarliestInUniverse\\
2026-09-11 & Permit content & the within-cell time gap between linked and unlinked permits, read at face value & most is immortal time: measured from the hearing, about \pcLmSurviveDays\% survives in means and \pcLmMedSurvive\% in medians\\
2026-09-11 & Data acquisition & the DR clock starts at the Planning record & at the DBI filing, a median \pcDrDiffMed\ days earlier\\
2026-09-11 & Exactions (in its own build) & the 2019 register as cited & the archive served the 2020 file for the 2019 capture; the fetcher now refuses another capture\\
2026-09-11 & The brief & about 1{,}040 pre-2010 items with a motion number and 590 without; 50{,}366 addresses & \chRrWith\ and \chRrWithout; \ccProbed\ addresses\\
2026-09-11 & The brief & AB 1763 among the Housing Accountability Act amendments & its own title is a density-bonus bill\\
2026-09-11 & The brief & a 25-unit ordinance of July 2026 & none among the Board's 2026 ordinances through Ord.~\rtSfLastOrd\\
\bottomrule\end{longtable}}
\fi
}

\title{The Market for Housing Regulation:\\
\large what has been built, what it established, and what is open}
\author{Dan Post}
\date{September 2026}

\begin{document}
\maketitle

\begin{abstract}
\noindent
A chronicle, not a paper: one document to reconstruct the project's state from, front to back.
It lists every memo with what later work revised in it, walks through them in the order they
were finished, inventories the data, states the model the project is building toward, maps each
model object to the data that would identify it, and keeps the register of what is open.

\medskip\noindent
\textbf{\chMemos\ memos and \chPagesTotal\ pages rest on one dataset of \chItems\ decisions.}
Every item the Planning Commission heard from \chItemFirst\ to \chItemLast\ with a case number is
extracted into structured fields (a model at \chHaikuTest\% field accuracy on a frozen test half),
dated to its hearing, and followed into DBI's permit record, the Planning Department's records,
the Commission's own motions and conditions, the parcel map, prices, fees and the law.

\medskip\noindent
\textbf{Several headline claims did not survive later work, and the timeline records each.} The
model case: ``no amount of work recovers condition content before 2010'' was a claim about two
URL patterns; a census of every case found text for \condReachEarlyHtml\% and
\condReachLateHtml\% of the conditioned items of the two pre-2010 eras --- little, not none.
Others: the two linkage routes do overlap; two of the project's own join keys were wrong; most of
the apparent slowness of Commission-linked permits is immortal time; the developer's clock on
the DR track starts at the DBI filing.

\medskip\noindent
\textbf{$\tau$ is now measured in dollars; $\sigma$ is measured as dispersion, not yet as the
model's primitive.} A new building of ten or more units owes a median \exPuMed\ per unit in
fees and inclusionary obligation; the developer's clock has a coefficient of variation of
\ckTotCvCu\ on the conditional-use track and \ckTotCvMin\ off the Commission. What blocks a first
draft is identification: separating $\tau$ from $\sigma$, the deterred projects, and treating
$\sigma$ as a parcel-level primitive (\S\ref{sec:chregister}).
\end{abstract}

\clearpage
\tableofcontents
\clearpage

% ═══════════════════════════════════════════════════════════════════════════
\section{The question}\label{sec:chquestion}

\textbf{The market for housing regulation.} A developer who wants to build in San Francisco
buys, in effect, a permission from the city, and the price has two parts. One is deterministic
and written down: fees, the inclusionary obligation, the template conditions every approval
carries. The other is stochastic: how long the process takes and what the Commission will ask
for. The project's premise is that the second part matters for what gets built, and where.

\textbf{The real-options frame.} A parcel is an option to develop. Uncertainty over delay and
conditions raises the value of waiting, so it raises the developer's filing threshold; and
because the cost of waiting is paid on the project's whole value, the threshold rises by
different amounts on different parcels. Variance in the regulatory process therefore does two
things a deterministic cost of the same expected size does not: it deters filing, and it
re-routes the filings that happen toward higher-value parcels.

\textbf{How the evidence is assembled.} The Commission's minutes, scraped for
\chItemFirst--\chItemLast, are parsed into agenda-item blocks; a deterministic stage dates each
block to its hearing and derives the sitting's attributes; a language model with retrieved
examples extracts each item's fields. From the item, four routes lead outward: to DBI's building
permits (a permit number printed in the minutes, the Planning records' own permit field, or a
parcel-and-date rule), to the Planning Department's records (case number), to the parcel map,
assessor roll, prices and zoning (block and lot), and to the Commission's own motions, where the
conditions of approval are written. The law --- fees, the inclusionary program, state statutes
--- enters as quoted claims checked against their primary sources. Three standards run through
all of it: extraction rules are frozen and hashed before they are validated, and in-sample and
out-of-sample figures are reported apart; a machine--label disagreement is checked against the
source before the machine is called wrong; and no number in a memo is typed by hand.

\textbf{Two papers.} Paper~1 estimates the within-jurisdiction mechanism in San Francisco,
where the Commission's record goes back to \chItemFirst\ and can be followed item by item. Paper~2
takes the primitives estimated there to the fragmentation question across the Bay Area's
jurisdictions. Everything in this chronicle serves Paper~1.

% ═══════════════════════════════════════════════════════════════════════════
\section{The memos, in order}\label{sec:chtimeline}

Table~\ref{tab:chtimeline} lists every memo by the date it was first completed, the question
it was asked, its answer in a sentence, and what later work revised or overturned in it. The
dates are those git first recorded each memo (for the three written in this batch, the run
log's); every memo has since been revised in place, and the revisions column is the point of the
table.

\chtable{timeline}

\subsection*{Every supersession, collected}

Table~\ref{tab:chsupersessions} collects every place where later work replaced a claim or a
figure, including the brief's own premises that this batch checked. Most are corrections of the
project's own instruments --- keys, parsers, scoring --- found when a number was checked against
its source; the rest are claims that generalised beyond what had been measured.

\chtable{supersessions}

% ═══════════════════════════════════════════════════════════════════════════
\section{A walk through them}\label{sec:chwalk}

Each subsection says what the memo asked, what it established, what it did not, and what it
changed downstream. Detail is in the memos.

\subsection{Meeting-level information (\chDateMeetingLevelInfo)}\label{sec:chmeetinglevelinfo}

\textit{Asked.} The item needs two facts it does not carry: the day it was heard, and the
character of the sitting. Both belong to the meeting. \textit{Established.} A deterministic date
stage that never re-derives block boundaries and was right on every gold item in four frozen
rounds; attribute fields in the mid-90s out of sample; a meeting table of \chMeetings\ sittings
over \chMeetFirst--\chMeetLast. Scored today against all \chMeetScored\ hand-confirmed meetings,
the attributes agree on \chMeetAcc\% of \chMeetValues\ values, \chMeetEarly\% for 1998--2001 and
\chMeetLate\% from 2015. \textit{Not established.} The all-meetings figure is unadjudicated, and
spot checks find hand labels that are wrong; two commissioners named Lee are left unresolved
rather than guessed. \textit{Downstream.} The date is the join key to every other panel, and the hearing sequence
of a case --- first hearing, continuances, final action --- is built from it; the meeting table
also carries who sat, which the corpus memo's commissioner analysis uses. The method's rule that
a corpus hole is drawn as a gap rather than interpolated became a standing rule of the project.

\subsection{Item-level extraction (\chDateExtractionMethodComparison)}\label{sec:chextractionmethodcomparison}

\textit{Asked.} Rules or a language model for the item fields? \textit{Established.} On a frozen
test half of a \chGoldItems-item gold set, a small model with retrieved examples scores
\chHaikuTest\% against the regular expressions' \chRegexTest\%; adjudicating the disagreements
found that about half were label errors, so an accuracy measured against unadjudicated gold is a
lower bound. \textit{Not established.} The modern era is thin in the gold set.
\textit{Downstream.} The corpus pass that produced the item table, and the rule that segmentation
stays deterministic while field semantics go to the model. Its storage layer (raw, interim and
clean replies kept apart) is what later made a silent data loss recoverable. Two of its figures
were corrected in the code review that followed (Table~\ref{tab:chsupersessions}).

\subsection{The corpus memo (\chDateDiscretionaryReviewPatterns)}\label{sec:chdiscretionaryreviewpatterns}

\textit{Asked.} What does the item-level dataset contain? \textit{Established.} Across
\chItems\ items and \corpusUniverse\ cases: outright denial has nearly vanished while conditioning
has more than tripled --- the Commission substituted a conditioning technology for a denial
technology; repeat appearances are a large share of the docket and the upper tail of time from
first to last hearing lengthened; discretionary review's share of the docket collapsed after the
2013 reforms, as that memo dates them. \textit{Not established.} Any ex-ante measure of burden: the data records realised
outcomes. \textit{Downstream.} It set the agenda: conditions (what is in them), linkage (what
happens after), delay (how much is forecastable). Its closing section named four routes from
realised outcomes to an ex-ante measure of burden; the real-options frame the later memos adopt
is the one that took. Its distinct-parcel count is superseded and not yet regenerated.

\subsection{Permit linkage (\chDatePermitLinkage)}\label{sec:chpermitlinkage}

\textit{Asked.} How far can a decision be followed into DBI's permits? \textit{Established.} The
printed permit number matches on \chPermitMatch\% of all \chPermitsDistinct\ distinct numbers the
minutes print; the apparent collapse of the route after 2019 was a change in how the minutes
print the number, not in the data. The route reaches discretionary reviews and almost no
conditional uses; a parcel-and-date fallback narrows the candidates at \chParcelPrecision\
precision and \chParcelRecall\ recall. Across DBI's rows, \chDbiCondTalk\ mention a Commission
condition at all. \textit{Not established.} The fallback's precision outside discretionary review,
where it was measured. \textit{Downstream.} The permit content memo keeps the fallback (T3) apart
from the printed number (T1) and the Planning records' permit field (T2), and calls only the first
two ``Commission-linked''. The negative result --- conditions are not in DBI --- sent the
conditions work to the Commission's own documents. The check that caught the post-2019
parsing failure, comparing the rate at which a thing is mentioned with the rate at which it is
parsed, became a standard diagnostic.

\subsection{Conditions of approval (\chDateConditionsOfApproval)}\label{sec:chconditionsofapproval}

\textit{Asked.} The item table records that a project was conditioned, not how. Where is the
content? \textit{Established.} Not in the minutes, not in DataSF, not in DBI; in the Commission
packet's draft motion, whose Exhibit~A numbers and names each condition, addressable by case
number from 2011. The conditions are templated, and a thirteen-category scheme reaches most
impositions; staff's conditions mostly change a project's obligations, the Commission's
amendments mostly change the building. \textit{Overturned.} Two claims. ``The two routes never
cover the same item'' was about what the minutes print; the Planning records' permit field reaches
\condBridgeCUPct\% of conditional-use items. And ``no amount of work recovers condition content
before 2010'' generalised from two URL patterns: the census found \condReachEarlyHtml\% and
\condReachLateHtml\% of the pre-2010 eras' conditioned items with text, against \condReachEarlyPdf\%
for 2011--2014, and \condResPreTenNone\ of the \condResPreTen\ pre-2010 items without text are
documents found at no address. \textit{Downstream.} The census, which replaced its sample with every case; its taxonomy
remains unvalidated against hand coding.

\subsection{Predicting delay (\chDatePredictingDelay)}\label{sec:chpredictingdelay}

\textit{Asked.} How much of the time from a case's first hearing to its last can be forecast from
what a developer sees on the day of the first? \textit{Established.} The level somewhat, the
tail barely, and in-sample fit roughly halves out of sample; the staff recommendation is the most
informative signal; the residual dispersion is large, and it is the first proxy for the cost of
uncertainty. \textit{Not established.} Anything causal. \textit{Superseded.} Its clock starts at
the first hearing; the data-acquisition memo's filing clocks and the clocks memo's distributions
replace it as the measure of duration.

\subsection{Data acquisition (\chDateDataAcquisition)}\label{sec:chdataacquisition}

\textit{Asked.} Where are the denominator, the clock, the scale, the price and the fees the
real-options frame needs? \textit{Established.} A defined risk set of \acqRiskParcels\ parcels and
\acqRiskParcelYears\ parcel-years carrying \acqRiskEvents\ hearings (\acqRiskRate\ per ten
thousand); a parcel join at \acqParcelJoin\% and a case join at \acqCaseAny\%; two filing clocks
that differ by a median of \acqClockGapMed\ days; a price on every parcel-year (out-of-sample
RMSE \acqPriceRmse\ in logs where imputed); fees that vary across sections within a year
(\acqFeeWithin\% of the log-rate variance) rather than over years (\acqFeeAbsYear\%).
\textit{Revised by itself.} Two conclusions reversed within days: the lettered-lot key and the
two-digit year had been diagnosed as data problems and were our own keys. \textit{Revised later.}
On the DR track the developer's clock starts at the DBI filing of the permit under review, a
median \pcDrDiffMed\ days before the Planning record (a dated note, 2026-09-11). \textit{Extended.}
The feasible envelope (\S\ref{sec:chbatch}).

\subsection{DataSF catalogue (\chDateDatasfPlanningCatalogue)}\label{sec:chdatasfplanningcatalogue}

\textit{Asked.} What does the portal hold under ``planning'' that the pipeline does not use?
\textit{Established.} \catTotal\ assets; no dataset of Commission actions, published or not;
\catPossibleNotCached\ assets flagged as possibly useful and not cached, of which the
Planning Department's record-level review-time metrics and DBI's station routing are delay data
nobody has measured against. \textit{Downstream.} None used yet; the review-time metrics are in the next steps.

\subsection{The conditions census (\chDateConditionsContent)}\label{sec:chconditionscontent}

\textit{Asked.} Replace the sample of twenty cases a year with every case and motion, and read
what the conditions say. \textit{Established.} \ccProbed\ addresses probed, \ccProbeHits\
documents found; condition text for \ccCondWithTextPct\% of the \ccCondItems\ conditioned items,
the earliest heard in \ccEarliestItemYear. The conditions are a template with a thin
project-specific layer, and what varies is numbers. The entitlement's term is written into the
template: a \ccValidityModal-month validity on \ccValidityModalShare\% of the motions that state one.
The parser was validated in frozen rounds; the third (this batch) scores \ccValInPrec\% precision
and \ccValInRecall\% recall in sample and \ccValOutPrec\% out of sample on rows heard
\ccValOutYearFrom--\ccValOutYearTo. \textit{Revised in this batch.} The document once read as the
earliest motion for a case in the universe was two motions bound together (Motion
\ccFixOldMotion, \ccFixOldDate, and a later one); every year-span resting on it moved.
\textit{Downstream.} The records request, the numbers table, and the exactions memo's
cross-check of the inclusionary schedule against the percentages the conditions print.

\subsection{Permit content (\chDatePermitsContent)}\label{sec:chpermitscontent}

\textit{Asked.} What is in DBI's table, and do the projects that reach the Commission differ from
similar ones? \textit{Established.} \pcUniverseShare\% of permits are development-relevant, and
\pcLinkedShare\% of those are Commission-linked. Linked projects are bigger and slower, but a
permit linked through a hearing held while it was pending had to survive to the hearing; measured
from the hearing on permits alive at it, about \pcLmSurviveDays\% of the within-cell gap in time
survives in means and \pcLmMedSurvive\% in medians. The lower abandonment of linked permits does
survive (\pcLinkedExit\% against \pcUnlinkedExit\% expired or withdrawn). \textit{Revised in this
batch.} The landmark comparison replaced a claim that almost none of the gap survives; cost
placeholders were dropped; the first-construction-document date was dropped as an outcome for
coverage.

\subsection{Clocks (\chDateClocks)}\label{sec:chclocks}

\textit{Asked.} The full distribution of every duration, by track, parcel value, size, period and
neighbourhood, with censoring handled. \textit{Established.} Kaplan--Meier medians of the
developer's whole clock of \ckTotMedCu\ days (conditional use), \ckTotMedDr\ (DR), \ckTotMedLpa\
(large projects) and \ckTotMedMin\ (ministerial); about a third of spells record no completion ten
years on (\ckTotStenPctCu\% on the conditional-use track); the coefficient of variation does not
fall consistently with parcel value. The outcome chain (added the same day): of housing projects
entitled \ocMatureFrom--\ocMatureTo, \ocDoneCu\% of conditional uses and \ocDoneMin\% of ministerial
permits record a completion. \textit{Not established.} $\sigma$ as the model defines it: a
cross-sectional CV mixes ex-ante uncertainty with what the developer knew.

\subsection{Exactions (\chDateExactions)}\label{sec:chexactions}

\textit{Asked.} What inclusionary requirement and impact fees vested for each project?
\textit{Established.} A version chain from 2002, every rate and vesting rule a quotation re-checked
against its source (\exClaims\ claims); a median obligation of \exPuMed\ per unit for new
buildings of ten or more units; bunching just below every inclusionary threshold (\exNineA\
buildings of nine units against \exTenA\ of ten since Proposition~C). \textit{Not established.}
The schedule agrees with the percentage the Commission's conditions print in \exCrossAgree\% of
comparable cases; area-specific rates and the grandfathering deadline's application are open.

\subsection{The regulatory timeline (\chDateRegulatoryTimeline)}\label{sec:chregulatorytimeline}

\textit{Asked.} Which laws generate variation this data can measure? \textit{Established.}
\rtLaws\ laws inventoried with quoted citations; \rtUsable\ usable, \rtNotMeasurable\ not
measurable here; the opt-in pathways are elections, not assignments. The docket's residential
conditional uses fell from \rtResCuPre\ a quarter to \rtResCuPost\ after the Housing Production
ordinances, so docket series need an era split there. \textit{Not established.} Any effect.

\subsection{The rest of this batch}\label{sec:chbatch}

\paragraph{The feasible envelope} (a section of the data-acquisition memo). \textit{Asked.} The
risk set models non-filing, and ``could have filed'' needs a capacity; the counterfactual needs
a benchmark that can be relaxed. \textit{Established.} For every parcel-year: the height limit
and bulk district read from the height district's name, the lot area read at the map lot so that
a condominium's unit parcels share their building's lot, the density rule of the zoning district
in force, the rear yard, and a maximum unit count, with the state density bonus apart. The rules
are quotations from the zoning control tables and the Government Code, checked like the
exaction claims. Under central assumptions the read districts allow \envUnitsTotal\ units in
\envYearLast, and \envSensLo--\envSensHi\ across the sensitivity runs. \textit{Not established.}
Rules before and after the reading of the tables it rests on; the named neighbourhood commercial
districts. \textit{Downstream.} The risk set's capacity and the counterfactual's benchmark.

\paragraph{The numbers inside the conditions} (a section of the conditions census).
\textit{Asked.} The census found that conditions are a template and that numbers are what vary;
extract them. \textit{Established.} \chRowsNumeric\ conditions with typed fields --- validity
periods, operating hours, bicycle and parking counts, stated inclusionary percentages and unit
counts, notice radii and periods, monitoring intervals, screening dimensions --- and their
distributions over time. \textit{Downstream.} The exactions memo's cross-check of its schedule
against the percentages the Commission's conditions print.

\paragraph{The outcome chain} (a section of the clocks memo). \textit{Asked.} Conditional on
entitlement, what share of approved projects and units is built, and how long does it take?
\textit{Established.} Completion shares by track and parcel value for projects entitled
\ocMatureFrom--\ocMatureTo; cumulative incidence of completion with exit as a competing event.
\textit{Not established.} Units delivered are a lower bound: a completion recorded under another of
a project's permit numbers is missed.

\paragraph{The records request.} \textit{Asked.} Turn the census's pre-2010 residual into a
request. \textit{Produced.} A draft letter to the Planning Department for the \chRrWith\
conditioned items heard before 2010 whose motion number is known and the \chRrWithout\ whose is
not, with a statement of everything already searched (\ccProbed\ addresses), and a separate set
of questions for the Assessor-Recorder on recorded notices of special restrictions. Nothing has
been sent.

% ═══════════════════════════════════════════════════════════════════════════
\section{What the data say so far}\label{sec:chfindings}

The substantive findings, collected from the memos that own them, each stated as the memo states
it and each a description.

\begin{itemize}
\item \textbf{The Commission conditions rather than denies.} Outright denial has nearly vanished
over the corpus while the share of items approved with conditions more than tripled (the corpus
memo); \ccCondWithTextPct\% of the \ccCondItems\ conditioned items now have their condition text
(the census).
\item \textbf{The conditions are a template.} Most imposed conditions carry one of a few dozen
headings; what varies is numbers. The entitlement's term is part of the template: a
\ccValidityModal-month validity on \ccValidityModalShare\% of the motions that state one.
\item \textbf{Delay is long, dispersed and hard to forecast.} From first filing to completion the
Kaplan--Meier median is \ckTotMedCu\ days on the conditional-use track and \ckTotMedMin\ off the
Commission, with coefficients of variation of \ckTotCvCu\ and \ckTotCvMin\ (the clocks memo);
what a developer observes at the first hearing forecasts the level of hearing delay somewhat and
its tail barely (the delay memo). About a third of spells record no completion ten years on.
\item \textbf{Commission projects differ from similar ones mostly in what they survive.} Linked
permits are bigger and slower, but most of the time difference is immortal time; their lower
abandonment (\pcLinkedExit\% against \pcUnlinkedExit\%) survives the correction (the permit
content memo).
\item \textbf{Most entitled housing is built.} Of housing projects entitled
\ocMatureFrom--\ocMatureTo, \ocDoneCu\% of conditional uses, \ocDoneDr\% of discretionary
reviews and \ocDoneLpa\% of large project authorisations record a completion; \ocExitCu\%,
\ocExitDr\% and \ocExitLpa\% exited (the outcome chain).
\item \textbf{The written price is large and sits at thresholds.} A new building of ten or more
units owes a median \exPuMed\ per unit, \exShareMed\% of DBI's valuation, mostly the
inclusionary fee; new buildings are fewer just above each inclusionary threshold than just
below it (the exactions memo).
\item \textbf{The docket changed character after the Housing Production ordinances.} Residential
conditional uses fell from \rtResCuPre\ to \rtResCuPost\ a quarter; SB~330's five-hearing cap does
not show in the hearing counts (\rtOverFivePre\% of residential cases heard more than five times
before, \rtOverFivePost\% after) (the regulatory timeline).
\item \textbf{Fees vary across the city, not over time.} A year effect absorbs \acqFeeAbsYear\% of the
variance in log fee rates; the variation is across sections and size bands (the acquisition memo).
\end{itemize}

% ═══════════════════════════════════════════════════════════════════════════
\section{The data inventory}\label{sec:chinventory}

Table~\ref{tab:chinventory} lists every table and panel the project has built or cached, its
unit of observation and size, its years, its known defects and the memo that documents it.
Everything lives under \texttt{\$MFHR\_DATA\_ROOT}, not in git.

\chtable{inventory}

% ═══════════════════════════════════════════════════════════════════════════
\section{The model as currently specified}\label{sec:chmodel}

The project is building toward a minimal model, deliberately narrower than the earlier
\texttt{dr\_supply\_model} blueprint (a document this repository does not hold, so its contents are
not summarised here):

\begin{quote}
One developer, one parcel, an option to file. The regulatory process imposes a deterministic cost
$\tau$ (fees, inclusionary obligations, the template conditions) and a stochastic delay and
outcome with variance $\sigma$. The exercise threshold depends on both, and the comparative
statics differ across parcel value: $\tau$ shifts the threshold roughly uniformly in project value,
$\sigma$ shifts it differentially. The counterfactual holds expected burden fixed and sets
$\sigma$ to zero.
\end{quote}

\paragraph{What each piece of data does.} In this model the data play four parts. The
exaction panel and the conditions census measure $\tau$: the dollars and the template obligations
a project owes, fixed at dates the law states. The clocks and the outcome chain measure the
distribution of delay and of outcomes that $\sigma$ summarises. The risk set, the price panel and
the envelope give the parcels that could have filed, what they were worth, and how much they
could hold --- the population on which the filing threshold is estimated. And the regulatory
timeline and the exaction panel's version chain supply the changes in the rules that might move
$\tau$ or $\sigma$ separately. The comparative statics the model delivers --- $\tau$ roughly
uniform in project value, $\sigma$ differential --- are what make the parcel-value gradient of
filing the central cross-sectional prediction, which is why every duration and obligation is
reported by parcel value.

\paragraph{Deferred.} To Paper~2 and to later drafts: strategic neighbours, endogenous supply of
discretionary review, the Commission's best response, and endogenous zoning. None of these is in
the model above, and none of the data built so far is organised around them.

\paragraph{An assumption, not a finding.} $\sigma$ is currently treated as a parcel-level
primitive estimated from a hazard. That is an assumption to be relaxed. The clocks memo measures
cross-sectional dispersion within cells; whether that dispersion is uncertainty the developer
faced, or heterogeneity the developer knew about, is not yet separated.

% ═══════════════════════════════════════════════════════════════════════════
\section{Which data object identifies which model object}\label{sec:chident}

Table~\ref{tab:chident} maps each primitive to the moment that would identify it, the dataset
that supplies the moment, and whether that dataset exists. The two rows where the answer is not
yet ``yes'' are the ones the open register calls blocking.

\chtable{ident}

% ═══════════════════════════════════════════════════════════════════════════
\section{The open register}\label{sec:chregister}

Every unresolved item, with what it needs --- data, a records request, a modelling decision or a
literature check --- and whether it blocks a first presentable draft.

\chtable{register}

\paragraph{The three that block.} \emph{Deterred projects}: the risk set models non-filing on
parcels with capacity (the envelope now supplies it), but the parcels whose owners would have
filed absent the regulatory process are not observed, and the filing hazard is estimated on
the parcels that did. \emph{Separating $\tau$ from $\sigma$}: the exactions memo and the
regulatory timeline name variation in $\tau$ with $\sigma$ plausibly fixed (the inclusionary
changes and their filing-date cutoffs) and variation in $\sigma$ (SB~330's hearing cap, the
Housing Production ordinances); none is yet an estimate. \emph{$\sigma$ as a primitive}: as
\S\ref{sec:chmodel} says.

% ═══════════════════════════════════════════════════════════════════════════
\section{Next steps}\label{sec:chnext}

In priority order, with a rough sense of scale.
\begin{enumerate}
\item \textbf{Write down the estimator for the filing hazard} on the risk set, with $\tau$ from the
exaction panel, the envelope as capacity, and parcel value from the price panel; decide how
$\sigma$ enters (a cell-level dispersion, or a latent parcel-level primitive). \emph{Weeks.}
\item \textbf{Turn the inclusionary cutoffs into a first design}: the notch and the filing-date
cutoffs, with the envelope held fixed. \emph{Days to a first cut; weeks to an estimate.}
\item \textbf{Join the Priority Equity Geographies polygon} and cut the Housing Production
ordinances' treated set. \emph{Days.}
\item \textbf{Validate T3 linkage outside discretionary review} on a hand-checked sample. \emph{Days.}
\item \textbf{Send the records request} (after editing) and the Assessor-Recorder questions.
\emph{Days to send; months to answer.}
\item \textbf{Close the exaction panel's gaps}: area-specific rates, TSF grandfathering,
existing-use credits, the 2017 register. \emph{Days each.}
\item \textbf{Use DataSF's review-time metrics and DBI's station routing} for the delay inside
the permit. \emph{Weeks.}
\end{enumerate}

% ═══════════════════════════════════════════════════════════════════════════
\section{Established and not established}\label{sec:chest}

\paragraph{Established.}
\begin{enumerate}
\item An item-level record of the Commission's decisions, \chItemFirst--\chItemLast, dated and
followed into permits, records, motions, parcels, prices, fees and law.
\item A dollar measure of $\tau$ from \exRegFirst\ and a distributional measure of the durations
on every track.
\item Which earlier claims were revised, and how (Table~\ref{tab:chtimeline}).
\end{enumerate}

\paragraph{Not established.}
\begin{enumerate}
\item $\sigma$ as the model defines it.
\item Any separation of $\tau$ from $\sigma$, and any correction for deterred projects.
\item Condition content before 2010 beyond the thin share now read.
\end{enumerate}

% ═══════════════════════════════════════════════════════════════════════════
\section{Reproducing this}\label{sec:chrepro}

\begin{verbatim}
cd code/commission_minutes_processing
python build_chronicle.py     # macros and tables for this memo
\end{verbatim}
It reads the other memos' generated macros files and tables, so each memo's own pipeline runs
first (each memo's ``Reproducing this'' section gives its commands); the dates come from git and
\texttt{run\_log.md}, the page counts from the compiled PDFs and the row counts from the data.

\end{document}
